% concatenated from WASEP-Xue-Zhao.tex
\documentclass[reqno,11pt]{amsart}
\usepackage{graphicx} % Required for inserting images
\usepackage{amsfonts, amscd, epsfig, amsmath, amssymb,enumerate,bm}
\usepackage{color}
\usepackage{mathrsfs}
\usepackage[margin=3cm]{geometry}
\usepackage[colorlinks=true,linkcolor=blue,citecolor=magenta]{hyperref}
\usepackage[displaymath, mathlines,pagewise]{lineno}
\usepackage{setspace}
\setstretch{1.1}
%\usepackage[notcite,notref]{showkeys}
\numberwithin{equation}{section}

\newtheorem{theorem}{Theorem}[section]
\newtheorem{lemma}[theorem]{Lemma}
\newtheorem{proposition}[theorem]{Proposition}
\newtheorem{corollary}[theorem]{Corollary}
\newtheorem{remark}[theorem]{Remark}
\newtheorem{definition}{Definition}[section]
\newtheorem{cond}{Condition}[section]
\newtheorem{assumption}{Assumption}[section]

\newcommand{\E}{\mathbb{E}}
\newcommand{\Pbb}{\mathbb{P}}
\newcommand{\R}{\mathbb{R}}
\newcommand{\Z}{\mathbb{Z}}

\newcommand{\Mcal}{\mathcal{M}}
\newcommand{\Scal}{\mathcal{S}}
\newcommand{\Ycal}{\mathcal{Y}}

\title[Moderate deviation principles]{Moderate deviation principles for the current and the tagged particle in the WASEP}
\keywords{Moderate deviation; tagged particle; current; exclusion process.}

\author{Xiaofeng Xue}
\address{School of Mathematics and Statistics, Beijing Jiaotong University, Beijing 100044, China.}
\email{xfxue@bjtu.edu.cn}

\author{Linjie Zhao}
\address{School of Mathematics and Statistics \& Hubei Key Laboratory of Engineering Modeling and Scientific Computing, Huazhong University of Science and Technology, Wuhan 430074, China.}
\email{linjie\_zhao@hust.edu.cn}

\thanks{\textbf{Acknowledgments.}  The project is supported by the National Natural Science Foundation of China with grant numbers 12401168 and 12371142, and the Fundamental Research Funds for the Central Universities in China.}


\begin{document}
	
	\begin{abstract}
		We study the weakly asymmetric simple exclusion process in one dimension. We prove  sample path moderate deviation principles for the current and the tagged particle  when the process starts from one of its stationary measures. We simplify the proof in our previous works \cite{xue2023moderate,xue2024sample}, where the same problem was investigated in the symmetric simple exclusion process. 
	\end{abstract}



\maketitle

\section{Introduction}

The exclusion process is one of the most studied interacting particle systems \cite{liggettips,liggett99}. In the dynamics, particles perform random walks on some graph subjected to the exclusion rule, that is, there is at most one particle at each site and jumping to occupied sites is suppressed.  It has been a long standing problem to study the behavior of a typical particle, usually called the tagged particle in the literature. The main challenge is that the tagged particle itself is not Markovian. By studying the environment process of the tagged particle, which is Markovian, much progress has been made, such as  law of large numbers \cite{saada1987tagged,rezakhanlou1994evolution}, central limit  theorems in equilibrium \cite{kipnis1986central,varadhan1995self,Kipnis86,sethuraman2000diffusive,sethuraman2006diffusive,zhao2023motion} and in nonequilibrium \cite{jara2008quenched,jara2006nonequilibrium,jara2009nonequilibrium}. Recently, large deviations \cite{sethuraman2013large,sethuraman2023atypical} and moderate deviations \cite{xue2023moderate,xue2024sample} for the current and the tagged particle in one dimension were also investigated. We refer the readers to \cite{komorowski2012fluctuations} and the above references for an excellent review on the related literature.

The aim of this paper is to simplify the proof of our previous works on sample path moderate deviation principles for the current and the tagged particle in the symmetric simple exclusion process in one dimension \cite{xue2023moderate,xue2024sample}.  To make our results more general, we study the weakly asymmetric simple exclusion process in one dimension. We show that when the process starts from one of its stationary measures, the current and the tagged particle process satisfy the moderate deviation principles under the correct time scaling.

The idea of the proof is as follows. Since particles cannot take over each other in the model, one can relate the current and the tagged particle to the empirical measure of the process. In \cite{zhao2024moderate}, the second author proved moderate deviation principles for the empirical measure of the process. Then, by contraction principle, the current and the tagged particle also satisfy moderate deviation principles. However, the rate functions for the current and the tagged particle are expressed as a  variational formula through this approach. By long calculations, the same authors solved it by using the Fourier approach in \cite{xue2023moderate,xue2024sample}. In this paper, we present a different approach, mainly relying on the contraction principle, and simplifies the previous proof significantly. The main observation is that the rate functions for the density fluctuation fields and for the limit of the density fluctuation fields are the same, which allows us to use the contraction principle back and forth. It seems that the approach developed here can also simplify the proof in \cite{zhao2024moderate}, where sample path moderate deviation principles for the occupation time of the exclusion process were proved in one dimension. 

The paper is organized as follows. In Section \ref{sec: model}, we introduce the process and the main results. We state moderate deviation principles for the empirical measure of the process in Section \ref{sec: density flieds}. Finally, the proof for moderate deviation principles of the current and the tagged particle is presented in Section \ref{sec: proof}.

\section{Model and Results}\label{sec: model}

We study the weakly asymmetric simple exclusion process (WASEP) on the one-dimensional integer lattice $\Z$. The state space of the process is $\Omega:= \{0,1\}^\Z$. For a configuration $\eta \in \Omega$, $\eta_x = 1$ if and only if there is one particle at site $x \in \Z$. Let $\alpha ,\beta,\gamma\geq 0$ be three parameters. Throughout the article, we shall take
\[\gamma = \min \{1+\beta, 2\}.\]
The infinitesimal generator of the WASEP is given by 
\[L_n = n^{\gamma} (L_s + \alpha n^{-\beta}L_a).\] 
Above, $L_s$ is associated to the dynamics of the symmetric simple exclusion process (SSEP), which acts on local functions $f: \Omega \rightarrow \R$ as
\[L_s f (\eta) = \frac{1}{2} \sum_{x \in \Z} \big[ f(\eta^{x,x+1} )- f(\eta)\big],\]
and $L_a$ is associated to the dynamics of the totally asymmetric simple exclusion process (TASEP),
\[L_a f (\eta) = \sum_{x \in \Z} \eta_x (1-\eta_{x+1})\big[ f(\eta^{x,x+1}) - f(\eta)\big],\]
where $\eta^{x,y}$ is the configuration obtained from $\eta$ by swapping the values of $\eta_x$ and $\eta_y$,
\[(\eta^{x,y})_{z} = \begin{cases}
 \eta_x, \quad \text{if } z = y,\\
  \eta_y, \quad \text{if } z = x,\\
   \eta_z, \quad \text{if } z \neq x, y.\\
\end{cases}\]
Note that if $\alpha = 0$ and $\gamma = 2$, then we obtain the  SSEP  in one dimension under the diffusive scaling; if $\beta = 0$ and $\gamma=1$, the we get the asymmetric simple exclusion process (ASEP) in one dimension under the hyperbolic scaling.



For $\rho \in [0,1]$, let $\nu_{\rho}$ be the product measure on the configuration space $\Omega$ with particle density $\rho$,
\[\nu_\rho (\eta_x = 1) = \rho, \quad  x \in \Z.\]
It is well known that the measure $\nu_\rho$ is reversible for the generator $L_s$ and is invariant for $L_a$, see \cite{liggettips} for example.


For any probability measure $\mu$ on $\Omega$, denote by $\Pbb_{\mu} \equiv \Pbb_{\mu}^n$ the probability measure on the path space $D(\R_+,\Omega)$ endowed with the Skorokhod topology corresponding to the law of the process $\eta(t)$ with generator $L_n$ and with initial distribution $\mu$. Let $\E_\mu \equiv \E_\mu^n$ be the corresponding expectation. 

We are interested in the current of the WASEP when it starts from its invariant measure $\nu_\rho$. For $x \in \Z$, the current $J^n_{x,x+1} (t)$ is defined as the number of particles jumping from $x$ to $x+1$ up to time $t$ minus the number of particles jumping from $x+1$ to $x$ up to time $t$. 

We shall also investigate the long-time behavior of the tagged particle. To define it, let $\nu_\rho^*$ be the measure $\nu_\rho$ conditioned on having a particle at the origin,
\[\nu_\rho^* (\cdot) = \nu_\rho (\cdot | \eta_0 = 1).\]
We call the particle initially at the origin the tagged particle. Under $\nu_\rho^*$, let $X^n (t)$ be the position of the tagged particle at time $t$. It is also well known that the measure $\nu_\rho^*$  is invariant for the environment process, defined as $\eta_{x + X^n (t)}(t)$ for $x \in \Z$,  as seen from the tagged particle, see \cite{liggettips} for example.

\subsection{Invariance principles} In this subsection, we state invariance principles for the current and the tagged particle.  In the rest of this article, we fix a time horizon $T > 0$. For any $t,s>0$, define the variance function as
\begin{equation}\label{a t s}
    a (t,s) = \begin{cases}
        \chi(\rho)\alpha|1-2\rho|\min\{t,s\} \quad &\text{if }\; \beta < 1,\\
        \chi(\rho)\left(f(t)+f(s)-f(|t-s|)\right) \quad &\text{if }\; \beta = 1,\\
        \chi(\rho)(\sqrt{t}+\sqrt{s}-\sqrt{|t-s|}) / \sqrt{2\pi}  \quad &\text{if }\; \beta > 1,
    \end{cases}
\end{equation}
 where $\chi (\rho) = \rho (1-\rho)$, and 
\[
f(t)=\frac{\alpha|1-2\rho|t}{2}+\mathbb{E}\left[\left(B_t-\alpha|1-2\rho|t\right)_+\right].
\]
Here, $\{B_t\}_{t\geq 0}$ is the standard one dimensional Brownian motion starting from the origin. For $r \in \R$, $r_+ := \max \{r,0\}$.  Note that $a (t,s) = 0$ when $\rho = 1/2$ and $\beta < 1$. Define
\[\bar{J}^n_{-1,0} (t) := J^n_{-1,0} (t) - t \alpha n^{\gamma-\beta}\chi (\rho), \quad \bar{X}^n (t) = X^n (t) - t \alpha n^{\gamma-\beta} (1-\rho).\]

By following \cite{gonccalves2008central} line by line, where fluctuations for the current and the tagged particle in the ASEP were investigated, one can prove the following result directly. The main idea is to relate the current and the tagged particle to the density fluctuation field of the process. For this reason, we omit the proof here.

\begin{proposition}\label{prop clt current tagged particle} Assume that $\rho \neq 1/2$ or $\beta \geq 1$. Then, as $n \rightarrow \infty$,  the sequence of processes $\{\bar{J}^n_{-1,0} (t)/\sqrt{n}, 0 \leq t \leq T \}_{n \geq 1}$ under $\mathbb{P}_{\nu_\rho}$ (respectively $\{\bar{X}^n (t)/\sqrt{n}, 0 \leq t \leq T \}_{n \geq 1}$ under $\mathbb{P}_{\nu_\rho^*}$) converges in the space $D([0,T],\R)$ to a Gaussian process with covariance $a(t,s)$ (respectively $a(t,s)/\rho^2$).
\end{proposition}

\begin{remark}
Note that when $\beta < 1$ and $\rho \neq 1/2$, the limiting process  is a  Brownian motion; when $\beta > 1$, the limit is a  fractional Brownian motion with Hurst parameter $1/4$.
\end{remark}

\begin{remark}
If $\beta=0$ and $\gamma = 1$, that is, for the ASEP, it was shown in \cite{balazs2010order} that  the variance of the current has order $n^{2/3}$ if $\rho = 1/2$.  For the TASEP, the one-point limit was proved in \cite{Johansson00} to be the Tracy-Widom distribution.
\end{remark}

\subsection{Moderate deviations}  In this subsection, we study moderate deviations for the current and the tagged particle.  
Let $\{a_n\}_{n \geq 1}$ be a sequence of real numbers such that \[\sqrt{n \log n} \ll a_n \ll n.\]  
Let $\mathbf{0}$ be the trajectory equals to zero at any time $0 \leq t \leq T$. Then, as $n \rightarrow \infty$,
\[\{\bar{J}^n_{-1,0} (t)/a_n, 0 \leq t \leq T \} \Rightarrow \mathbf{0}, \quad \{\bar{X}^n (t)/a_n, 0 \leq t \leq T \} \Rightarrow \mathbf{0}.\]
Moderate deviations concern the rate of the above convergence.

To introduce the MDP rate functions for the current and the tagged particle, we first recall the following  representation for $\{B_t^{1/4}, t \geq 0\}$ the fractional Brownian motion with Hurst parameter $1/4$ (see \cite{Decreu1999FBM} for example),
\[B_t^{1/4} = \int_0^t \mathcal{K} (t,s) d B_s, \quad t \geq 0,\]
where the kernel $ \mathcal{K} (t,s)$ is defined as
\[
\mathcal{K}(t,s)=\frac{(t-s)^{-1/4}}{\sqrt{V}\Gamma(3/4)}F(1/4, -1/4, 3/4, 1-\tfrac{t}{s}), \quad 0\leq s<t.
\]
Above, 
\[
V=\frac{8\Gamma(3/2)\cos(\pi/4)}{\pi}, \quad 
F(\alpha,\beta,\gamma,z)=\sum_{k=0}^{+\infty}\frac{(\alpha)_k(\beta)_k}{(\gamma)_kk!}z^k
\]
with $(a)_k := \Gamma(a+k)/\Gamma(a)$. 

If $\beta < 1$, we define $\mathcal{H}_\beta = C^1 ([0,T])$, the space of continuously differentiable functions on $[0,T]$. For $h \in \mathcal{H}_\beta$, denote by $\dot{h}$ the usual derivative of $h$.  In this case, the rate function for the current is defined as
	\[\mathcal{J}^\beta (h):= 
\frac{1}{2 \chi(\rho) \alpha |1-2\rho|}  \|\dot{h}\|_{L^2 ([0,T])}^2 \]
if $h \in \mathcal{H}_\beta$, and $= + \infty$ otherwise.

If $\beta > 1$, define $\mathcal{H}_\beta$ as the family of functions $h: [0,T] \rightarrow \R$ satisfying that there exists $\dot{h} \in L^2 ([0,T])$ such that
\[h (t) = \int_0^t \mathcal{K}(t,s) \dot{h} (s) ds, \quad \forall 0 \leq t \leq T.\]
In this case, the rate function for the current is defined as
	\[\mathcal{J}^\beta (h):= 	\frac{\sqrt{\pi}}{2  \sqrt{2} \chi(\rho)}  \|\dot{h}\|_{L^2 ([0,T])}^2.\]
if $h \in \mathcal{H}_\beta$, and $= + \infty$ otherwise. 

Finally, the rate function for the tagged particle is defined as
 \[\mathcal{X}^\beta (\cdot) := \rho^2 \mathcal{J}^\beta (\cdot).\]



Now, we state the main result of this article. 

\begin{theorem}\label{thm mdp current tagged particle} Assume that $\beta < 1, \rho \neq 1/2$ or $\beta > 1$. Then, the sequence  of processes $\{\bar{J}^n_{-1,0} (t)/a_n, 0 \leq t \leq T \}_{n \geq 1}$ under $\mathbb{P}_{\nu_{\rho}}$  (respectively $\{\bar{X}^n (t)/a_n, 0 \leq t \leq T \}_{n \geq 1}$ under $\mathbb{P}_{\nu_{\rho}^*}$) satisfies the moderate deviation principles in the space $D([0,T], \R)$ with decay rate $a_n^2/n$ and with rate function $\mathcal{J}^\beta$ (respectively with rate function $\mathcal{X}^\beta$). 	
	
Precisely speaking, for any closed set $\mathcal{C} \subset D([0,T], \R)$ and for any open set $\mathcal{O} \subset D([0,T],\R)$,
	\begin{align*}
		\limsup_{n \rightarrow \infty} \frac{n}{a_n^2} \log \mathbb{P}_{\nu_{\rho}} \Big(\{\bar{J}^n_{-1,0} (t)/a_n, 0 \leq t \leq T \} \in \mathcal{C}\Big) \leq - \inf_{h \in \mathcal{C}} \mathcal{J}^\beta (h),\\
		\liminf_{n \rightarrow \infty} \frac{n}{a_n^2} \log \mathbb{P}_{\nu_{\rho}} \Big(\{\bar{J}^n_{-1,0} (t)/a_n, 0 \leq t \leq T \} \in \mathcal{O}\Big) \geq - \inf_{h \in \mathcal{O}} \mathcal{J}^\beta (h).
	\end{align*}

	The precise meaning of the moderate deviation principles for the tagged particle is similar.
\end{theorem}


\begin{remark}\label{rmk: beta equals 1}
For $\beta = 1$,  the moderate deviation principles still hold for the current and the tagged particle.  The decay rate is $a_n^2/n$ as before. The rate functions are implicitly given by the following variational formulas: for $h \in D([0,T],\R)$,
\begin{align*}
	\mathcal{J}^\beta (h) &=\inf \Big\{\frac{1}{2} \mathbf{h}^T A^{-1} \mathbf{h}:  m \geq 1, 0 \leq t_1 < t_2 < \ldots < t_m \leq T, t_1,\ldots,t_m \in \Delta_c (h)\Big\},\\
	\mathcal{X}^\beta (h) &= \rho^2 \mathcal{J}^\beta (h),
\end{align*}
where $\mathbf{h} = (h(t_1), \ldots, h(t_m))^T \in \R^m$, $A = (a(t_i,t_j))_{1 \leq i,j \leq m}$ with $a(\cdot,\cdot)$ defined in \eqref{a t s}, and $\Delta_c (h)$ is the set of continuous points of $h$.  However, we are not aware of how to solve the above infimum explicitly so far.
\end{remark}

\begin{remark}
The above theorem should hold in the regime $\sqrt{n} \ll a_n \ll n$.  We need the technical assumption $a_n \gg \sqrt{n \log n}$ in order to prove the exponential tightness of the current and the tagged particle processes. However, we are not aware of how to remove this technical assumption.
\end{remark}

%
%\hrule
%
%This section removed
%
%\section{Proof of Theorem \ref{thm clt current tagged particle}}\label{section 3}
%
%
%
%\subsection{Invariance principle for the density fluctuation fields} To prove Theorem \ref{thm clt current tagged particle}, we need to study the  density fluctuation fields of the WASEP, which acts on Schwartz functions $H \in \Scal (\R)$  as
%\[\Ycal^n_t (H) = \frac{1}{\sqrt{n}} \sum_{x \in \Z} \Bar{\eta}_x (t)H(\tfrac{x}{n}),\]
%where $\Bar{\eta}_x = \eta_x - \rho$.
%
%The following result concerns the stationary fluctuations for the WASEP.  Its proof is standard in the theory of hydrodynamic limits, see \cite[Chapter 11]{klscaling}  for example.  For this reason, we omit the proof here.
%
%\begin{theorem}\label{thm clt filed}
%Under $\Pbb_{\nu_\rho}$, the sequence of processes $\{\Ycal^n_t, 0 \leq t \leq T\}_{n \geq 1}$ converges in distribution in the space $D([0,T], \Scal^\prime (\R))$, as $n \rightarrow \infty$, to the process $\{\Ycal_t, 0 \leq t \leq T\}$, which is the unique solution to the following SPDE:
%\begin{enumerate}[(i)]
%    \item if $\beta > 1$, then
%     \[\partial_t \Ycal_t = \frac{1}{2} \Delta \Ycal_t + \sqrt{\chi (\rho)} \nabla d \mathcal{W}_t;\]
%     \item if $\beta = 1$, then
%     \[\partial_t \Ycal_t = \frac{1}{2} \Delta \Ycal_t - \alpha(1-2\rho) \nabla \Ycal_t+ \sqrt{\chi (\rho)} \nabla d \mathcal{W}_t;\]
%     \item if $\beta < 1$, then
%     \[\partial_t \Ycal_t  = - \alpha(1-2\rho) \nabla \Ycal_t\]
%\end{enumerate}
%with initial distribution $\mathcal{Y}_0$ such that the distribution of $\mathcal{Y}_0 (H)$ is normal  with mean $0$ and variance $\chi(\rho)\langle H, H\rangle$ for any $H\in \Scal(\R)$. 
%\end{theorem}
%
%\begin{remark}\label{remark covariance}
%By solving the $\Scal^\prime(\R)$-valued SPDEs in Theorem \ref{thm clt filed}, we have the following covariance formulas.  For any $G, H\in \Scal(\R)$,  any $t \geq s$,
%\begin{enumerate}[(i)]
%    \item if  $\beta<1$, then
%\begin{equation}\label{equ covariance sub}
% {\rm Cov}(\mathcal{Y}_t(H), \mathcal{Y}_s(G))=\chi(\rho) \langle S_{t-s}H, G\rangle, 
%\end{equation}
%where
%\[
%S_tH(u)=H(u+\alpha(1-2\rho)t), \quad u\in \mathbb{R}, \;t\geq 0;
%\]
%\item  if $\beta>1$, then
%\begin{equation}\label{equ covariance sup}
% {\rm Cov}(\mathcal{Y}_t(H), \mathcal{Y}_s(G))=\chi(\rho)\langle T_{t-s}H, G\rangle, 
%\end{equation}
%where
%\[
%T_tH(u)=\mathbb{E}[H(u+B_t)], \quad u\in \mathbb{R}, \;t\geq 0,
%\]
%with $\{B_t\}_{t\geq 0}$ being the standard Brownian motion starting from the origin;
%\item  if $\beta=1$, then
%\begin{equation}\label{equ covariance cri}
% {\rm Cov}(\mathcal{Y}_t(H), \mathcal{Y}_s(G))=\chi(\rho)\langle S_{t-s}T_{t-s}H, G\rangle.
%\end{equation}
%\end{enumerate}
%\end{remark}
%
%Calculations of the above covariance functions are  given in Appendix \ref{Appendix A.1}.
%
%\begin{remark}\label{remark extend to L2}
%    Let $\{\mathcal{Y}_t\}_{t\geq 0}$ be  given in Theorem \ref{thm clt filed}. Later, we will also let $\mathcal{Y}_t$ act on $L^2(\mathbb{R})$ functions, which is well defined since the process is stationary. Indeed, for any $H\in L^2(\mathbb{R})$, let $\{H_n\}_{n\geq 1}$ be a sequence in $\Scal(\R)$ such that $H_n\rightarrow H$ in $L^2(\R)$. According to Equations \eqref{equ covariance sub}-\eqref{equ covariance cri}, in each case $\{\mathcal{Y}_t(H_n)\}_{n\geq 1}$ is a $L^2(\mathbb{P_{\nu_\rho}})$-Cauchy sequence. Hence, it is natural to define $\mathcal{Y}_t(H)$ as the $L^2(\mathbb{P_{\nu_\rho}})$-limit of $\{\mathcal{Y}_t(H_n)\}_{n\geq 1}$.
%\end{remark}
%
%For any integer $l > 0$, define
%\begin{equation}\label{equ 3.4}
%G_l(u)=(1-\tfrac{u}{l})\mathbf{1}_{\{0\leq u\leq l\}}, \quad u\in \mathbb{R}.
%\end{equation}
%In the next subsection, we need to consider $\Ycal_t (\mathbf{1}_{\{u \geq 0\}})$. However, the above quantity makes no sense since  the indicator function $\mathbf{1}_{\{u \geq 0\}}$ does not belong to $L^2 (\R)$. It turns out that $G_l$ is a good approximation to the above indicator function, which has been used in many other contexts, see \cite{jara2006nonequilibrium,gonccalves2008central} for example.
%
%The next results calculate the limit of the covariacne functions of the density fluctuation fileds when acting on the approximating functions $G_l$ as $l \rightarrow +\infty$.
%
%\begin{lemma}\label{lem current variance}
%Let $\{\mathcal{Y}_t\}_{t\geq 0}$ be the fluctuation limit defined in Theorem \ref{thm clt filed}. Then, for any $s,t > 0$,
%\[\lim_{l\rightarrow+\infty}{\rm Cov}\left(\mathcal{Y}_t(G_l)-\mathcal{Y}_0(G_l), \mathcal{Y}_s(G_l)-\mathcal{Y}_0(G_l)\right) = a(t,s),\]
%where $a(\cdot,\cdot)$ was defined in \eqref{a t s}.
%\end{lemma}
%
%
%\begin{proof}[Proof of Lemma \ref{lem current variance}]
%We only check the case $\beta = 1$ and the remaining cases can be proved similarly. Without loss of generality, we assume that $t>s$ and $1-2\rho>0$. By Equation \eqref{equ covariance cri},
%  \begin{multline}\label{equ 3.7}
%      {\rm Cov}\left(\mathcal{Y}_t(G_l)-\mathcal{Y}_0(G_l), \mathcal{Y}_s(G_l)-\mathcal{Y}_0(G_l)\right)\\
%      =\chi(\rho)\left(\langle G_l-S_tT_tG_l, G_l\rangle+\langle G_l-S_sT_sG_l, G_l\rangle-\langle G_l-S_{t-s}T_{t-s}G_l, G_l\rangle\right). 
%  \end{multline}
%We calculate the first term inside the above parentheses,
%\begin{align*}
%\langle G_l-S_tT_tG_l, G_l\rangle&=\int_0^l \Big(1-\frac{u}{l} \Big)\left( \Big(1-\frac{u}{l}\Big)-\mathbb{E}\big[G_l(u+B_t+\alpha(1-2\rho)t)\big]\right)du\\
%&={\rm \uppercase\expandafter{\romannumeral1}}+{\rm \uppercase\expandafter{\romannumeral2}},
%\end{align*}
%where 
%\begin{align*}
%    {\rm \uppercase\expandafter{\romannumeral1}} &=\int_0^l \Big(1-\frac{u}{l} \Big)^2\mathbb{P}\left(u+B_t+\alpha(1-2\rho)t\leq 0 \text{~or~}\geq l\right)du,\\
%    {\rm \uppercase\expandafter{\romannumeral2}}&=\int_0^l \Big(1-\frac{u}{l} \Big)\mathbb{E}\left[ \frac{B_t+\alpha(1-2\rho)t}{l} \mathbf{1}_{\{0\leq u+B_t+\alpha(1-2\rho)t\leq l\}}\right] du. 
%\end{align*}
%By the dominated convergence theorem,
%\begin{align}\label{equ 3.8}
%\lim_{l\rightarrow+\infty}{\rm \uppercase\expandafter{\romannumeral1}}&=\int_0^{+\infty}\mathbb{P}(u+B_t+\alpha(1-2\rho)t\leq 0)du \notag\\
%&=\int_0^{+\infty}\mathbb{P}(B_t\geq u+\alpha(1-2\rho)t)du\notag\\
%%&=\mathbb{E}\left\{\left(\int_0^{B_t-\alpha(1-2\rho)t}1du\right)1_{\{B_t\geq \alpha(1-2\rho)t\}}\right\} \notag\\
%&=\mathbb{E}\left[\left(B_t-\alpha(1-2\rho)t\right)_+\right],
%\end{align}
%where $r_+ = \max \{r,0\}$ for $r \in \R$. We decompose ${\rm \uppercase\expandafter{\romannumeral2}}$ as ${\rm \uppercase\expandafter{\romannumeral3}}+{\rm \uppercase\expandafter{\romannumeral4}}$, where 
%\begin{align*}
%    {\rm \uppercase\expandafter{\romannumeral3}} &=\mathbb{E} \Big[ (B_t+\alpha(1-2\rho)t) l^{-1}\int_0^{l-B_t-\alpha(1-2\rho)t} \Big(1-\frac{u}{l} \Big)du \mathbf{1}_{\{B_t+\alpha(1-2\rho)t\geq 0\}}\Big],\\
%    {\rm \uppercase\expandafter{\romannumeral4}} &=\mathbb{E} \Big[ (B_t+\alpha(1-2\rho)t) l^{-1} \int_{-B_t-\alpha(1-2\rho)t}^{l} \Big(1-\frac{u}{l} \Big)du \mathbf{1}_{\{B_t+\alpha(1-2\rho)t\leq 0\}}\Big].
%\end{align*}
%Using the dominated convergence theorem again,
%\begin{align*}
%\lim_{l\rightarrow+\infty}{\rm \uppercase\expandafter{\romannumeral3}}&=\lim_{l\rightarrow+\infty}\mathbb{E} \Big[ (B_t+\alpha(1-2\rho)t) \mathbf{1}_{\{B_t+\alpha(1-2\rho)t\geq 0\}}
%\frac{1}{2}\Big(1-\Big(\frac{B_t+\alpha(1-2\rho)t}{l}\Big)^2\Big)\Big]  \\
%&=\frac{1}{2}\mathbb{E}\big[(B_t+\alpha(1-2\rho)t)_+ \big].
%\end{align*}
%Similarly,
%\[
%\lim_{l\rightarrow+\infty}{\rm \uppercase\expandafter{\romannumeral4}}=-\frac{1}{2}\mathbb{E}\big[ (B_t+\alpha(1-2\rho)t)_- \big],
%\]
%where $r_- = -\min \{r,0\}$ for $r \in \R$. Therefore, 
%\[\lim_{l\rightarrow+\infty}{\rm \uppercase\expandafter{\romannumeral2}}=\frac{1}{2}\mathbb{E}\left[B_t+\alpha(1-2\rho)t\right]=\frac{\alpha(1-2\rho)t}{2},\] 
%and hence
%\begin{equation}\label{equ 3.9}
%   \lim_{l\rightarrow+\infty}\langle G_l-S_tT_tG_l, G_l\rangle=\mathbb{E}\left[\left(B_t-\alpha(1-2\rho)t\right)_+\right]+\frac{\alpha(1-2\rho)t}{2}.
%\end{equation}
% Consequently, the result follows from Equations \eqref{equ 3.7} and \eqref{equ 3.9}.
%\end{proof}
%
%\subsection{Invariance principles for the current}  In this subsection, we investigate invariance principles for the current in Theorem \ref{thm clt current tagged particle}. Since the number of particles is conserved,
%\begin{equation}
%    \eta_x (t) - \eta_x (0) = J^n_{x-1,x} (t) - J^n_{x,x+1} (t), \quad x \in \Z.
%\end{equation}
%Thus,
%\begin{equation}
%\begin{aligned}
%	\sum_{x} G_l (\tfrac{x}{n}) [\eta_x (t) - \eta_x (0)] &= \sum_{x} G_l (\tfrac{x}{n}) [J^n_{x-1,x} (t) - J^n_{x,x+1} (t)]\\
%	&= \sum_{x} [G_l (\tfrac{x+1}{n}) - G_l (\tfrac{x}{n})] J^n_{x,x+1} (t) \\
%	&= J^n_{-1,0} (t) - \frac{1}{nl} \sum_{x=0}^{nl-1} J^n_{x,x+1} (t).
%\end{aligned}
%\end{equation}
%Note that, for each $x$,
%\[M^n_{x,x+1} (t) = J^n_{x,x+1} (t) - \int_0^t \Big\{ \frac{n^\gamma}{2} (\eta_x - \eta_{x+1}) + \alpha n^{\gamma-\beta} \eta_x (1-\eta_{x+1}) \Big\} ds\]
%is a martingale with quadratic variation 
%\begin{equation}\label{qv martingale current}
%    \langle M^n_{x,x+1}\rangle(t) = \int_0^t \Big\{ \frac{n^\gamma}{2} (\eta_x - \eta_{x+1})^2 + \alpha n^{\gamma-\beta} \eta_x (1-\eta_{x+1}) \Big\} ds.
%\end{equation}
%Here, we simply write $\eta_x (s)$ as $\eta_x$ when there are no confusions. This permits us to write the current as
%\begin{align}\label{current expression}
%\Bar{J}_{-1,0}^n (t) = \sqrt{n}[\Ycal^n_t (G_l) - \Ycal^n_0 (G_l)]+ \frac{1}{nl} \sum_{x=0}^{nl-1} M^n_{x,x+1} (t) + I_n (t),
%\end{align}
%where $I_n (t) = I_{n,1} (t) + I_{n,2} (t)$ with
%\begin{align*}
%    I_{n,1} (t) &= \frac{n^{\gamma}}{2nl} \sum_{x=0}^{nl-1} \int_0^t (\eta_x - \eta_{x+1}) ds,\\
%    I_{n,2} (t) &= \frac{\alpha n^{\gamma-\beta}}{nl} \sum_{x=0}^{nl-1} \int_0^t [\eta_x(1-\eta_{x+1}) - \chi(\rho)] ds.
%\end{align*}
%
%The following result allows us to approximate the current by functionals of the density fluctuation fields.
%
%\begin{lemma}\label{lem: current approx} 
%\[\limsup_{l \rightarrow +\infty} \limsup_{n \rightarrow +\infty} \E_{\nu_\rho} \Big[ \sup_{0 \leq t \leq T} \Big( \frac{\Bar{J}^n_{-1,0}(t)}{\sqrt{n}} - [\Ycal^n_t (G_l) - \Ycal^n_0 (G_l)] \Big)^2 \Big] = 0.\]
%\end{lemma}
%
%\begin{proof} By \eqref{current expression} and Cauchy-Schwarz inequality, we only need to prove the estimates in the lemma for the last two terms on the right-hand side of \eqref{current expression}.
%	
%	 
%   We first deal with the martingale term in \eqref{current expression}. Since there are no common jumps between the martingales $M^n_{x,x+1} (t)$, they are independent. Then, by Doob's inequality and \eqref{qv martingale current}, 
% \begin{align*}
%     \E_{\nu_\rho} \Big[ \sup_{0 \leq t \leq T} \Big( \frac{1}{n^{3/2}l} \sum_{x=0}^{nl-1} M^n_{x,x+1} (t) \Big)^2 \Big] \leq  \frac{C}{n^3l^2} \sum_{x=0}^{nl-1}   \E_{\nu_\rho} \Big[ \Big(  M^n_{x,x+1} (T) \Big)^2 \Big] 
%     \leq \frac{CTn^{\gamma-2}}{l},
% \end{align*}
% which converges to zero as $n \rightarrow +\infty, l \rightarrow +\infty$ since $\gamma \leq 2$.
%
% For $I_{n,2} (t)$, we use Cauchy-Schwarz inequality, and bound 
% \begin{align*}
%     \E_{\nu_\rho} \Big[ \sup_{0 \leq t \leq T} \big(I_{n,2}(t) / \sqrt{n} \big)^2 \Big] &\leq \frac{Cn^{2\gamma-2\beta}T^2}{n^3l^2}  E_{\nu_\rho} \Big[ \Big(\sum_{x=0}^{nl-1} \{\eta_x(1-\eta_{x+1}) - \chi(\rho)\} \Big)^2 \Big]\\
%     &\leq \frac{CT^2n^{2\gamma-2\beta-2}}{l}. 
% \end{align*}
% The last term also vanishes in the limit since $\gamma \leq 1 + \beta$.
%
% It remains to bound the term $I_{n,1} (t)$. By Kipnis-Varadhan inequality \cite[Lemma 2.4, Section 2.5]{komorowski2012fluctuations}, 
% \begin{equation*}
%     \E_{\nu_\rho} \Big[  \sup_{0 \leq t \leq T} \Big( I_{n,1} (t) / \sqrt{n} \Big)^2 \Big] \leq \frac{CTn^{2\gamma}}{n^3l^2} \sup_{f \in L^2 (\nu_\rho)} \Big\{ \sum_{x=0}^{nl-1}  E_{\nu_\rho} \Big[ (\eta_x - \eta_{x+1}) f(\eta)   \Big] + \langle L_nf,f\rangle_{\nu_\rho} \Big\}.
% \end{equation*}
% Since $\nu_\rho$ is an invariant measure for the dynamics,
% \begin{align*}
%     \langle - L_nf,f\rangle_{\nu_\rho}  \geq n^{\gamma} \langle - L_s f,f\rangle_{\nu_\rho} = \frac{n^\gamma}{4} \sum_{x \in \Z} E_{\nu_\rho} \Big[ (f(\eta^{x,x+1}) - f(\eta))^2\Big] =: n^\gamma D_s (f).
% \end{align*}
% By making the change of transformations $\eta \mapsto \eta^{x,x+1}$, and by Young's inequality, for any $A > 0$, 
% \begin{align*}
%     \sum_{x=0}^{nl-1}  E_{\nu_\rho} \Big[ (\eta_x - \eta_{x+1}) f(\eta) \Big] &= \frac{1}{2} \sum_{x=0}^{nl-1}  E_{\nu_\rho} \Big[ (\eta_x - \eta_{x+1}) \big( f(\eta) - f(\eta^{x,x+1}) \big) \Big]\\
%     &\leq \sum_{x=0}^{nl-1} \Big\{ \frac{A}{4} E_{\nu_\rho} \Big[ (f(\eta^{x,x+1}) - f(\eta))^2\Big] + \frac{1}{4A}  \Big\}\\
%     &\leq A D_s (f) + \frac{nl}{4A}.
% \end{align*}
% Thus,
% \[     \E_{\nu_\rho} \Big[  \sup_{0 \leq t \leq T} \Big( I_{n,1} (t) / \sqrt{n} \Big)^2 \Big] \leq \frac{CTn^{2\gamma}}{n^3l^2} \sup_{f \in L^2 (\nu_\rho)} \Big\{ A D_s (f) + \frac{nl}{4A} - n^\gamma D_s (f) \Big\}.\]
% Take $A=n^\gamma$, then
% \begin{align*}
%     \E_{\nu_\rho} \Big[ \sup_{0 \leq t \leq T}  \Big( I_{n,1} (t) / \sqrt{n} \Big)^2 \Big] \leq \frac{CTn^{\gamma-2}}{l}.
% \end{align*}
% Since $\gamma \leq 2$, the proof is concluded. 
%\end{proof}
%
%To conclude the proof for the  invariance principle of the current, we only need to show that
%\begin{enumerate}[(i)]
%	\item the sequence of processes $\{\Bar{J}^n_{-1,0} (t)/\sqrt{n}, 0 \leq t \leq T\}_{n \geq 1}$ is tight in $D([0,T],\R)$;
%	\item the finite-dimenional distribution of $\{\Bar{J}^n_{-1,0} (t)/\sqrt{n}, 0 \leq t \leq T\}$ converges to a Gaussian vector with covariance  $a(t,s)$ defined in \eqref{a t s}.
%\end{enumerate}
%Statement (i)  follows from Lemma \ref{lem: current approx} and the tightness of the density fluctuation field.  For Statement (ii), by using Lemma \ref{lem: current approx} and Theorem \ref{thm clt filed},  the finite-dimenional distribution of $\{\Bar{J}^n_{-1,0} (t)/\sqrt{n}, 0 \leq t \leq T\}$ converges to the limit of the corresponding finite-dimensional distribution of $\{\mathcal{Y}_t (G_l) - \mathcal{Y}_0 (G_l)\}$  as $n \rightarrow + \infty, l \rightarrow +\infty$. Then, we conclude the proof by Lemma \ref{lem current variance}. 
%
%\begin{remark}
%	By the basic coupling \cite{liggettips}, one can couple  two processes with initial conditions $\nu_{\rho}$ and $\nu_\rho^*$ such that they differ at only one site at any time $t \geq 0$. Thus, the results proved so far are also true under the measure $\Pbb_{\nu_\rho^*}$. See \cite[Remark 2.2]{xue2023moderate} for details.
%\end{remark}
%
%\subsection{Invariance principles for the tagged particle} As in the last subsection, we need to prove that under $\mathbb{P}_{\nu_{\rho}^*}$, 
%\begin{enumerate}[(i)]
%	\item the sequence of processes $\{\Bar{X}^n (t)/\sqrt{n}, 0 \leq t \leq T\}_{n \geq 1}$ is tight in $D([0,T],\R)$;
%	\item the finite-dimenional distribution of $\{X^n (t)/\sqrt{n}, 0 \leq t \leq T\}$ converges to a Gaussian vector with covariance  $a(t,s)/\rho^2$.
%\end{enumerate}
%
%We start with the following observation: for $x > 0$, 
%\begin{align}
% \{X^n (t) > x\} = \Big\{J^n_{-1,0} (t) \geq \sum_{y=0}^{x} \eta_y (t) \Big\},  \quad \{X^n (t) < -x\} = \Big\{J^n_{-1,0} (t) <  \sum_{y=-x}^{-1} \eta_{y} (t) \Big\}.
%\end{align}
%
%\hrule



\section{Density fluctuation fields}\label{sec: density flieds}

In this section, we study the density fluctuation fields of the process. The  density fluctuation field $\Ycal^n_t$ of the WASEP, which acts on Schwartz functions $H \in \Scal (\R)$, is defined  as
\[\Ycal^n_t (H) = \frac{1}{\sqrt{n}} \sum_{x \in \Z} \Bar{\eta}_x (t)H(\tfrac{x}{n}),\]
where $\Bar{\eta}_x = \eta_x - \rho$.

The following result concerns the stationary fluctuations for the WASEP.  Its proof is based on the  martingale approach, which is standard in the theory of hydrodynamic limits, see \cite[Chapter 11]{klscaling}  for example.  For this reason, we omit the proof here.

\begin{proposition}\label{prop clt filed}
	Under $\Pbb_{\nu_\rho}$, the sequence of processes $\{\Ycal^n_t, 0 \leq t \leq T\}_{n \geq 1}$ converges in distribution in the space $D([0,T], \Scal^\prime (\R))$, as $n \rightarrow \infty$, to the process $\{\Ycal_t, 0 \leq t \leq T\}$, which is the unique solution to the following SPDE:
	\begin{enumerate}[(i)]
		\item if $\beta > 1$, then
		\[\partial_t \Ycal_t = \frac{1}{2} \Delta \Ycal_t + \sqrt{\chi (\rho)} \nabla d \mathcal{W}_t;\]
		\item if $\beta = 1$, then
		\[\partial_t \Ycal_t = \frac{1}{2} \Delta \Ycal_t - \alpha(1-2\rho) \nabla \Ycal_t+ \sqrt{\chi (\rho)} \nabla d \mathcal{W}_t;\]
		\item if $\beta < 1$, then
		\[\partial_t \Ycal_t  = - \alpha(1-2\rho) \nabla \Ycal_t.\]
	\end{enumerate}
	Moreover, the  initial distribution $\mathcal{Y}_0$ of the above processes satisfies that  the distribution of $\mathcal{Y}_0 (H)$ is normal  with mean $0$ and variance $\chi(\rho)\langle H, H\rangle$ for any $H\in \Scal(\R)$. 
\end{proposition}

Next, we study large deviations for the limiting process $\mathcal{Y}_t$ in Proposition \ref{prop clt filed}. We first introduce the  rate function.  For  $\beta>0$ and $\mu\in D ([0, T], \Scal^\prime(\R))$, define
\begin{equation}\label{equ mdp rate function of the fluctuation field}
    \mathcal{Q}^\beta(\mu)=\mathcal{Q}_{\rm ini}(\mu_0)+\mathcal{Q}_{\rm dyn}^\beta(\mu),
\end{equation}
where $\mathcal{Q}_{\rm ini}$ corresponds to the deviation from the initial state, and $\mathcal{Q}_{\rm dyn}^\beta$ comes from the evolution of the dynamics. The initial rate function $\mathcal{Q}_{\rm ini}$ is defined as 
\[
\mathcal{Q}_{\rm ini}(\mu_0)=\sup_{h\in \Scal(\R)}\left\{\mu_0 (h)-\frac{\chi(\rho)}{2}\int_\mathbb{R}h^2(u)du\right\}, \quad \mu_0 \in \Scal^\prime(\R)
\]
To define $\mathcal{Q}_{\rm dyn}^\beta$, we introduce the following positive-definite quadratic form:  for given $T>0$ and $G,H\in D ([0, T], \Scal(\R))$,  define
\[
[G, H]:= \int_0^T\langle \nabla G_t, \nabla H_t \rangle dt,
\]
where $\langle \cdot,\cdot \rangle$ is the usual inner product in $L^2 (\R)$. Let $\mu\in D ([0, T], \Scal^\prime(\R))$.
\begin{itemize}
	\item For $\beta<1$,  define
	\begin{equation}\label{equ dynamic mdp rate sub}
		\mathcal{Q}_{\rm dyn}^\beta(\mu)=\sup_{H\in C_c^{1,+\infty}([0, T]\times \mathbb{R})}\left\{\mu_T(H_T)-\mu_0(H_0)-\int_0^T\mu_t\left((\partial_t+\alpha(1-2\rho)\nabla)H_t\right)dt\right\}.    
	\end{equation}
	\item For $\beta>1$, define
	\begin{equation}\label{equ dynamic mdp rate sup}
		\mathcal{Q}_{\rm dyn}^\beta(\mu)
		=\sup_{H\in C_c^{1,+\infty}([0, T]\times \mathbb{R})}\left\{\mu_T(H_T)-\mu_0(H_0)-\int_0^T\mu_t\left((\partial_t+\frac{1}{2}\Delta)H_t\right)dt-\frac{\chi(\rho)}{2}[H, H]\right\}.  
	\end{equation}
	\item For $\beta=1$, define
	\begin{equation}\label{equ dynamic mdp rate cri}
		\mathcal{Q}_{\rm dyn}^1(\mu)
		=\sup_{H\in C_c^{1,+\infty}([0, T]\times \mathbb{R})}\left\{\mu_T(H_T)-\mu_0(H_0)-\int_0^T\mu_t\left((\partial_t+\mathcal{P})H_t\right)dt-\frac{\chi(\rho)}{2}[H, H]\right\}, 
	\end{equation}
	where $\mathcal{P}=\frac{1}{2}\Delta+\alpha(1-2\rho)\nabla$. 
\end{itemize}
Below, we write $\mathcal{Q}_{\rm dyn}^\beta, \mathcal{Q}^\beta$ as $\mathcal{Q}_{{\rm dyn}, T}^\beta, \mathcal{Q}_T^\beta$ respectively when we need to emphasize the dependence of the rate functions on the time horizon $T$. 

For later use, we give some properties of the above rate functions.  For any $G,H \in C_c^{1, +\infty}([0, T]\times \R)$, define the equivalence relation  
\[G\sim H \quad \text{if and only if} \quad [G-H, G-H]=0.\] 
We denote by $\mathbb{H}$ the completion of $C_c^{1, +\infty}([0, T]\times \R)/\sim$ under the inner product $[\cdot, \cdot]$. 

\begin{lemma}\label{lemma finite rate} 
$(1)$ If $\mu_0 \in \Scal^\prime(\R)$ satisfies that $\mathcal{Q}_{\rm ini}(\mu_0)<+\infty$, then there exists $\varphi \in L^2(\R)$ such that
\[
\mu_0 (H) = \langle \varphi, H\rangle, \; \forall H \in L^2 (\R), \quad \text{and} \quad \mathcal{Q}_{\rm ini}(\mu_0)=\frac{1}{2\chi(\rho)} \|\varphi\|_{L^2 (\R)}^2.
\]
$(2)$ If $\mu\in D ([0, T], \Scal^\prime(\R))$ satisfies that $\mathcal{Q}^\beta(\mu)<+\infty$, then $\mu\in C([0, T], L^2(\R))$. Moreover, we have the following characterizations of  $\mu$ and the rate function.

\begin{itemize}
	\item If  $\beta>1$, then there exists $G\in \mathbb{H}$ such that $\mu$ is the unique weak solution to the PDE
	\[\begin{cases}
	\partial_t \mu(t,u)=\frac{1}{2}\Delta\mu(t,u)-\Delta G(t,u), \text{~}t> 0, u\in \mathbb{R},\\
	\mu(0,u) = \varphi (u), \; u \in \R,
	\end{cases}\]
where $\varphi$ is identified in $(1)$. Moreover, $\mathcal{Q}^\beta_{\rm dyn}(\mu)=[G,G] / (2 \chi(\rho))$.
\item If  $\beta=1$, then there exists $G\in \mathbb{H}$ such that $\mu$ is the unique weak solution to the PDE
\[\begin{cases}
	\partial_t \mu(t,u)=\mathcal{P}^* \mu(t,u)-\Delta G(t,u), \text{~}t\geq 0, u\in \mathbb{R},\\
	\mu(0,u) = \varphi (u), \; u \in \R.
\end{cases}\]
Here, $\mathcal{P}^*$ is the adjoint of $\mathcal{P}$ in $L^2 (\R)$. Moreover, $\mathcal{Q}^\beta_{\rm dyn}(\mu)=[G,G] / (2 \chi(\rho))$.
\item If  $\beta< 1$, then 
\[
\mu(t,u)=\varphi( u-\alpha(1-2\rho)t), \text{~}t\geq 0, u\in \mathbb{R}.
\]
Moreover, $\mathcal{Q}^\beta_{\rm dyn}(\mu)=0$.
\end{itemize}
\end{lemma}


\begin{proof}[Proof of Lemma \ref{lemma finite rate}]
We only prove the case $\beta<1$ in (2) since the remaining statements follow from  Riesz representation theorem directly, see \cite{gao2003moderate, kipnis1989hydrodynamics} for example. If there exists $H\in C_c^{1,+\infty}([0, T]\times \mathbb{R})$ such that
\[
\mu_T(H_T)-\mu_0(H_0)-\int_0^T\mu_t\left((\partial_t+\alpha(1-2\rho)\nabla)H_t\right)dt\neq 0,
\]
then
\[
\mu_T(aH_T)-\mu_0(aH_0)-\int_0^T\mu_t\left((\partial_t+\alpha(1-2\rho)\nabla)(aH_t)\right)dt\rightarrow+\infty
\]
as $a\rightarrow+\infty$ or $a\rightarrow -\infty$, and hence $\mathcal{Q}^\beta_{\rm dyn}(\mu)=+\infty$. Consequently, if $\mathcal{Q}^\beta(\mu)<+\infty$, we must have
\[
\mu_T(H_T)-\mu_0(H_0)-\int_0^T\mu_t\left((\partial_t+\alpha(1-2\rho)\nabla)H_t\right)dt=0
\]
for any $H\in C_c^{1,+\infty}([0, T]\times \mathbb{R})$ and hence $\mathcal{Q}_{\rm dyn}^\beta(\mu)=0$.

Take the test function $H$ with the form $H(t,u)=b_th(u)$ for some $h\in C_c^\infty(\R), b\in C^1 ([0, T])$, then
\[
b_T\mu_T(h)-b_0\mu_0(h)-\int_0^Tb^\prime_t\mu_t(h)dt=\int_0^Tb_t\alpha(1-2\rho)\mu_t(\nabla h)dt.
\]
Since $b$ is arbitrary, we have
\[
\frac{d}{dt}\mu_t(h)=\alpha(1-2\rho)\mu_t(\nabla h).
\]
For $s\in \R$, let $S_{s}$ be the translation operator:  $S_{s}h(u)=h(u+\alpha(1-2\rho)s)$.  Then, for any $t\geq 0$,
\[
\frac{d}{dt}\mu_t(S_{-t}h)=\alpha(1-2\rho)\mu_t(\nabla S_{-t}h)-\alpha(1-2\rho)\mu_t(\nabla S_{-t}h)=0.
\]
 Hence,
\[
\mu_t(h)=\mu_0 (S_th)=\int_{\R}\varphi(u-\alpha(1-2\rho)t)h(u)du
\]
for any $h\in C_c^\infty(\R)$, thus concluding the proof. 
\end{proof}



Since the process $\mathcal{Y}_t$ is Gaussian, the following result is straightforward and thus the proof is omitted.  

\begin{lemma}\label{lem MPD of the fluctuation limit}
The sequence of processes $\left\{\frac{\sqrt{n}}{a_n}\mathcal{Y}_t:~0\leq t\leq T\right\}_{n\geq 1}$ satisfies the large deviation principles with decay rate $a_n^2/n$ and with rate function $\mathcal{Q}^\beta$. 
\end{lemma}

Finally, we study moderate deviations for the  density fluctuation field. The rescaled density fluctuation field $\mu^n=\{\mu^n_t\}_{0\leq t\leq T}$ of the process is defined as  
\begin{equation}\label{equ intial rate funtion}
	\mu^n_t(du)=\frac{1}{a_n}\sum_{x\in \mathbb{Z}}\overline{\eta}_x(t)\delta_{x/n}(du)
\end{equation}
where $\delta_a(du)$ is the Kronecker Dirac measure concentrated at the point $a$.

\begin{proposition}\label{prop MPD of the scaled density field}
	The rescaled density field $\{\mu^n_t, 0 \leq t \leq T\}_{n\geq 1}$ satisfies the moderate deviation principles in the space $D([0,T],\mathcal{S}^\prime (\R))$ with decay rate $a_n^2 / n$ and with rate function $\mathcal{Q}^\beta$.    
\end{proposition}

 Its proof is similar to \cite{zhao2024moderate}, where the MDP was proved for the process with generator $n^2 (L_s + \alpha n^{-\beta} L_a)$. For that reason, we omit the proof here.



\section{Proof of Theorem \ref{thm mdp current tagged particle}}\label{sec: proof}

In this section, we prove moderate deviation principles for the current and the tagged particle. In Subsection \ref{subsec: current}, we prove sample path moderate deviation principles for the current by assuming the exponential tightness and finite dimensional moderate deviation principles of the current. The exponential tightness for the current are proved in Subsection \ref{subsec: pf lemma 4.1}, and the finite dimensional moderate deviation principles are proved in Subsection \ref{subsec:pf lemma 4.2}.  Finally, sample path moderate deviation principles for the tagged particle are outlined in Subsection \ref{subsec:tagged}.

\subsection{MDP for the current}\label{subsec: current} To prove the sample path MDP for the current, we only need to prove the exponential tightness and finite-dimensional MDP for the process $\{\bar{J}^n_{-1,0} (t)/a_n, 0 \leq t \leq T\}_{n \geq 1}$, which are summarized in the following two lemmas. 

\begin{lemma}\label{lem: exp tight current}
The sequence of processes $\{\bar{J}^n_{-1,0}(t) / a_n, 0 \leq t \leq T\}_{n \geq 1}$ is exponentially tight.
\end{lemma} 

\begin{lemma}\label{lem: finiteD mdp current}
For any $m \geq 1$ and for any $0 \leq t_1 < t_2 < \ldots < t_m \leq T$, the sequence of random vectors $\{\bar{J}^n_{-1,0} (t_i)/a_n, 1 \leq i \leq m\}_{n \geq 1}$ satisfies the MDP with decay rate $a_n^2 / n$ and with rate function $\mathcal{J}^{\beta,m} \equiv \mathcal{J}^\beta_{\{t_i\}_{i=1}^m}$, where  for $\mathbf{r} =(r_1,\ldots, r_m)^T\in \mathbb{R}^m$,
\[\mathcal{J}^{\beta,m}  (\mathbf{r}) = \frac{1}{2} \mathbf{r}^T A^{-1} \mathbf{r}.\]
Here, $A = A_{\{t_i\}_{i=1}^m} = (a(t_i,t_j))_{1 \leq i,j \leq m}$ is the $m \times m$ matrix with $a(\cdot,\cdot)$ defined in \eqref{a t s}. 
\end{lemma}

Now, we prove the sample path MDP for the current by using the above two lemmas.

\begin{proof}[Proof of Theorem \ref{thm mdp current tagged particle} for the current]
By Lemmas \ref{lem: exp tight current}, \ref{lem: finiteD mdp current} and \cite[Theorem 4.28]{Feng2006LDP}, the sequence of processes $\{\bar{J}^n_{-1,0} (t)/a_n, 0 \leq t \leq T\}_{n \geq 1}$ satisfies the MDP with decay rate $a_n^2/n$ and with rate function
\[\mathcal{J}^\beta (r(\cdot)):=\inf \Big\{\frac{1}{2} \mathbf{r}^T A^{-1} \mathbf{r}:  m \geq 1, 0 \leq t_1 < t_2 < \ldots < t_m \leq T, t_1,\ldots,t_m \in \Delta_c (r(\cdot))\Big\},\]
where $\Delta_c (r(\cdot))$ is the set of continuous points of the function $r: [0,T] \rightarrow \R$, $\mathbf{r} = (r(t_1),\ldots,r(t_m))^T \in \R^m$ and $A=  (a(t_i,t_j))_{1 \leq i,j \leq m}$. Moreover, since $a(t,s)$ is the covariance function of the Brownian motion if $\beta < 1$, and that of the fractional Brownian motion with Hurst parameter $1/4$ if $\beta > 1$, using \cite[Theorem 4.28]{Feng2006LDP} again, the last infimum equals a constant multiple the LDP rate function of the Brownian motion (respectively the fractional Brownian motion with Hurst parameter $1/4$ ) if $\beta < 1$ (respectively if $\beta > 1$). Precisely,
\[\mathcal{J}^\beta (r(\cdot)) = \begin{cases}
\frac{1}{2 \chi(\rho) \alpha |1-2\rho|}  \|\dot{r}\|_{L^2 ([0,T])}^2 \quad &\text{if}\; \beta < 1,\\
\frac{\sqrt{\pi}}{2  \sqrt{2} \chi(\rho)}  \|\dot{r}\|_{L^2 ([0,T])}^2 \quad &\text{if}\; \beta > 1.\\
\end{cases}\]
This concludes the proof.
\end{proof}



\subsection{Proof of Lemma \ref{lem: exp tight current}}\label{subsec: pf lemma 4.1} We follow the proof of  \cite[Lemma 5.2]{xue2024sample},  where the exponential tightness was proved for the current in the SSEP. The proof in the previous paper used the stirring representation for the SSEP, which does not hold any more for the WASEP.  Despite that, most of the proof follows \cite[Lemma 5.2]{xue2024sample} line by line. Thus, we only sketch it here and underline the main differences.


It suffices to prove the following two estimates:
\begin{itemize}
	\item \begin{equation}\label{exp tight c1}
		\lim_{M \rightarrow \infty} \limsup_{n \rightarrow \infty} \frac{n}{a_n^2} \log \Pbb_{\nu_\rho} \Big( \sup_{0 \leq t \leq T} |\bar{J}^n_{-1,0} (t) | > a_n M\Big) = -\infty;
	\end{equation}
\item for any $\varepsilon > 0$,
\begin{equation}\label{exp tight c2}
\lim_{\delta \rightarrow 0} \limsup_{n \rightarrow \infty} \sup_{\tau \in \mathcal{T}_T} \frac{n}{a_n^2} \log \Pbb_{\nu_\rho} \Big(\sup_{0 \leq t \leq \delta} |\bar{J}^n_{-1,0} (t+\tau) - \bar{J}^n_{-1,0} (\tau)| >a_n \varepsilon\Big) = - \infty,
\end{equation}
where $\mathcal{T}_T$ is the family of all stopping times bounded by $T$. 
\end{itemize}

For any integer $l > 0$, define
\begin{equation}\label{G l}
	G_l(u)=(1-\tfrac{u}{l})\mathbf{1}_{\{0\leq u\leq l\}}, \quad u\in \mathbb{R}.
\end{equation}
Since the number of particles is conserved,
\begin{equation*}
	\eta_x (t) - \eta_x (0) = J^n_{x-1,x} (t) - J^n_{x,x+1} (t), \quad x \in \Z.
\end{equation*}
Multiplying by $G_l (x/n)$ on both hand sides, then summing over $x \in \Z$ and using the summation by parts formula, we have
\begin{equation}\label{current density filed relation}
	\begin{aligned}
		\sum_{x} G_l (\tfrac{x}{n}) [\eta_x (t) - \eta_x (0)] &= \sum_{x} G_l (\tfrac{x}{n}) [J^n_{x-1,x} (t) - J^n_{x,x+1} (t)]\\
		&= \sum_{x} [G_l (\tfrac{x+1}{n}) - G_l (\tfrac{x}{n})] J^n_{x,x+1} (t) \\
		&= J^n_{-1,0} (t) - \frac{1}{nl} \sum_{x=0}^{nl-1} J^n_{x,x+1} (t).
	\end{aligned}
\end{equation}
Note that the expectation of both hand sides in the last equation is zero with respect to $\Pbb_{\nu_\rho}$. Thus,
\begin{equation}\label{current decom}
\bar{J}^n_{-1,0} (t) / a_n = \langle \mu^n_t,G_l\rangle -\langle \mu^n_0,G_l  \rangle + \frac{1}{na_nl} \sum_{x=0}^{nl-1} \bar{J}^n_{x,x+1} (t).
\end{equation}


By following the proof of \cite[Lemma 5.1]{xue2024sample}, one can show that for any $\varepsilon > 0$,
\begin{equation}\label{super exp 1}
\limsup_{l\rightarrow+\infty} \limsup_{n \rightarrow \infty} \frac{n}{a_n^2} \log \Pbb_{\nu_\rho} \Big( \sup_{0 \leq t \leq T} \Big| \frac{1}{na_n l} \sum_{x=0}^{nl-1} \bar{J}^n_{x,x+1} (t) \Big| > \varepsilon \Big) = - \infty.
\end{equation}
In particular, the third term on the right hand side of \eqref{current decom} is exponentially tight as $n \rightarrow \infty, l \rightarrow \infty$.

To conclude the proof, we only need to show that the estimates in \eqref{exp tight c1} and \eqref{exp tight c2} are true with $\bar{J}^n_{-1,0} (t)$ replaced by $\langle \mu^n_t,G_l\rangle$. The problem is that the function $G_l$ does not belong to $\mathcal{S} (\R)$.  Thus, we need to approximate it by Schwartz functions. It is in this step that we need the technical assumption $a_n \gg \sqrt{n \log n}$. Let $\tilde{G}_l \in  \mathcal{S} (\R)$ satisfying
\[\|\tilde{G}_l  - G_l\|_{L^2(\R)} \leq Cl^{-1}, \quad {\rm supp} (\tilde{G}_l) \subset [-2\ell,2\ell].\]
By Proposition \ref{prop MPD of the scaled density field}, for any $l$, $\langle \mu^n_t,\tilde{G}_l\rangle$ is exponentially tight. Thus, we only need to  prove that,  for any $\varepsilon > 0$, 
 \begin{equation}\label{exp tight error}
\limsup_{n \rightarrow \infty} \frac{n}{a_n^2} \log \Pbb_{\nu_\rho} \Big( \sup_{0 \leq t \leq T} |\langle \mu^n_t,F_l\rangle | >  \varepsilon \Big) = -\infty,
\end{equation}
where  $F_l := \tilde{G}_l - G_l$. 

Let $t_i = i T / n^3$ for $0 \leq i \leq n^3$. Then, we bound the probability in \eqref{exp tight error} by 
\[\Pbb_{\nu_\rho} \Big( \sup_{0 \leq i \leq n^3} |\langle \mu^n_{t_i},F_l\rangle | >  \varepsilon/2 \Big) + \Pbb_{\nu_\rho} \Big( \sup_{0 \leq t \leq T} |\langle \mu^n_t,F_l\rangle | - \sup_{0 \leq i \leq n^3} |\langle \mu^n_{t_i},F_l\rangle |  >  \varepsilon/2 \Big).\]
By using the assumption $a_n \gg \sqrt{n \log n}$, we have
\[\limsup_{n \rightarrow \infty} \frac{n}{a_n^2} \log \Pbb_{\nu_\rho} \Big( \sup_{0 \leq i \leq n^3} |\langle \mu^n_{t_i},F_l\rangle | >  \varepsilon/2 \Big)  = - \infty,\]
see \cite[Proof of the first term in (5.16)]{xue2024sample} for details.  

To deal with the second term above, we need to modify the proof in \cite{xue2024sample} slightly since we cannot use the stirring representation for the SSEP here. Since the process is stationary,
\begin{align*}
	&\Pbb_{\nu_\rho} \Big( \sup_{0 \leq t \leq T} |\langle \mu^n_t,F_l\rangle | - \sup_{0 \leq i \leq n^3} |\langle \mu^n_{t_i},F_l\rangle |  >  \varepsilon/2 \Big) \\
	\leq& \Pbb_{\nu_\rho} \Big( \sup_{0 \leq i \leq n^3-1} \,\sup_{t_i \leq t \leq t_{i+1}} |\langle \mu^n_t - \mu^n_{t_i},F_l\rangle |  >  \varepsilon/2 \Big)\\
	\leq& n^3 \Pbb_{\nu_\rho} \Big( \sup_{0 \leq t \leq Tn^{-3}} |\langle \mu^n_t - \mu^n_{0},F_l\rangle |  >  \varepsilon/2 \Big).
\end{align*}
By basic coupling (see \cite{liggettips} for example), we can assume that there is a particle at the origin initially. We label this particle by $Y_0$ and then label the other particles from the left to the right in an increasing order. Let $Y_i (t)$ be the position of the particle with label $i \in \Z$ at time $t$. Since particles cannot take over each other, $Y_i (t) < Y_{i+1} (t)$ for any $i \in \Z$ and any $t \geq 0$. We rewrite
\[\langle \mu^n_t - \mu^n_{0},F_l \rangle = \frac{1}{a_n} \sum_{i \in \Z} \Big\{F_l \big(\tfrac{Y_i(t)}{n}\big) - F_l \big(\tfrac{Y_i(0)}{n}\big)\Big\}.\]
Next, we consider the two cases $|i| > 3 n l$ and $|i| \leq 3 n l$ respectively. For the first case, let $\xi$ be a Poisson random variable with parameter $\alpha T n^{-1}$. Then, $\sup_{0 \leq t \leq Tn^{-3}}|Y_i (t) - Y_i (0)|$ is stochastically bounded by $\xi$.  Note that $F_l$ is supported in $[-2l,2l]$ and $|Y_i (0)| \geq 3 n l$ for $|i| > 3 n l$. Thus, $F_l (Y_i (0) / n) = 0$ for  $|i| > 3 n l$.  By standard large deviation estimates, 
\begin{align*}
	&\Pbb_{\nu_\rho} \Big( \sup_{0 \leq t \leq Tn^{-3}} \big|\frac{1}{a_n} \sum_{|i| > 3 n l} \Big\{F_l \big(\tfrac{Y_i(t)}{n}\big) - F_l \big(\tfrac{Y_i(0)}{n}\big)  \big|  >  \varepsilon/2 \Big)\\
	\leq& \sum_{|i| > 3 n l} \Pbb_{\nu_\rho} \Big(  \sup_{0 \leq t \leq Tn^{-3}} |Y_i (t) - Y_i (0)| > |i| - 2 n l\Big)\\
	\leq& \sum_{|i| > 3 n l} \Pbb_{\nu_\rho} (\xi > |i| - 2 n l) \leq C e^{-Cnl}
\end{align*}
for some constant $C > 0$. To deal with the sum over  $|i| \leq 3 n l$, we introduce the event
\[A_n = \Big\{  \sup_{|i| \leq 3 n l} \sup_{0 \leq t \leq Tn^{-3}} |Y_i (t) - Y_i (0)| \leq a_n^{3/2}/n^{1/2}\Big\}.\]
Then,
\[\Pbb_{\nu_\rho}  (A_n^c) \leq (3nl+1) \Pbb_{\nu_\rho} (\xi > a_n^{3/2}n^{-1/2}) \leq C n l e^{-Ca_n^{3/2}n^{-1/2}}.\]
We claim that for $n$ large enough, 
\[A_n \cap \Big\{   \sup_{0 \leq t \leq Tn^{-3}} \big|\frac{1}{a_n} \sum_{|i| \leq 3 n l} \Big\{F_l \big(\tfrac{Y_i(t)}{n}\big) - F_l \big(\tfrac{Y_i(0)}{n}\big)  \big|  >  \varepsilon/2 \Big\} = \emptyset.\]
Adding up the above estimates,  for $n$ large enough,
\[	\Pbb_{\nu_\rho} \Big( \sup_{0 \leq t \leq T} |\langle \mu^n_t,F_l\rangle | - \sup_{0 \leq i \leq n^3} |\langle \mu^n_{t_i},F_l\rangle |  >  \varepsilon/2 \Big) \leq C n^3 e^{-C n l} + C n^4 l e^{-Ca_n^{3/2}n^{-1/2}}.\]
Since $a_n \ll n$,
\[\lim_{n \rightarrow \infty} \frac{n}{a_n^2} \log \Pbb_{\nu_\rho} \Big( \sup_{0 \leq t \leq T} |\langle \mu^n_t,F_l\rangle | - \sup_{0 \leq i \leq n^3} |\langle \mu^n_{t_i},F_l\rangle |  >  \varepsilon/2 \Big) = - \infty.\]

It remains to prove the claim. For any $t$, let $B_{n,t}$ be the set of labels $i \in [-3nl,3nl]$ such that $Y_i (0) < 0, Y_i (t) \geq 0$ or  $Y_i (0) > 0, Y_i (t) \leq 0$.  Since the orderings of the particles are preserved by the dynamics, on the event $A_n$ we have $|B_{n,t}| \leq C a_n^{3/2} n^{-1/2}$ for any $0 \leq t \leq T n^{-3}$. Thus, on the event $A_n$,
\[\sup_{0 \leq t \leq Tn^{-3}} \big|\frac{1}{a_n} \sum_{|i| \leq 3 n l, i \in B_{n,t}} \Big\{F_l \big(\tfrac{Y_i(t)}{n}\big) - F_l \big(\tfrac{Y_i(0)}{n}\big)  \Big\} \big| \leq C(l) a_n^{1/2} n^{-1/2}.\]
Note that $F(l)$ is only discontinuous at the origin. For $i \notin B_{n,t}$, by the piecewise smoothness of $F_l$, on the event $A_n$,
\[\sup_{0 \leq t \leq Tn^{-3}} \big|\frac{1}{a_n} \sum_{|i| \leq 3 n l, i \notin B_{n,t}} \Big\{F_l \big(\tfrac{Y_i(t)}{n}\big) - F_l \big(\tfrac{Y_i(0)}{n}\big)  \Big\} \big| \leq C(l) a_n^{1/2} n^{-1/2}. \]
Thus, on the event $A_n$, 
\[\sup_{0 \leq t \leq Tn^{-3}} \big|\frac{1}{a_n} \sum_{|i| \leq 3 n l} \Big\{F_l \big(\tfrac{Y_i(t)}{n}\big) - F_l \big(\tfrac{Y_i(0)}{n}\big)  \Big\} \big| \leq C(l) a_n^{1/2} n^{-1/2},\]
which cannot be larger than $\varepsilon / 2$ for $n$ large enough since $a_n \ll n$. This proves the claim. 

\subsection{Proof of Lemma \ref{lem: finiteD mdp current}}\label{subsec:pf lemma 4.2} Intuitively, since 
\[\bar{J}^n_{-1,0} (t) / a_n = \frac{1}{a_n} \sum_{x=0}^{+\infty} [\eta_x (t) - \eta_x (0)] = \langle \mu^n_t - \mu^n_0, \chi_{[0,\infty)}\rangle,\]
by contraction principle, $\{\bar{J}^n_{-1,0} (t_i)/a_n, 1 \leq i \leq m\}$ should satisfy the MDP with decay rate $a_n^2/n$ and with rate function
\[\inf\left\{\mathcal{Q}_{t_m}^\beta(\mu):~ \langle \mu_{t_i} -\mu_{0},  \chi_{[0,\infty) } \rangle=r_i\text{~for all~}1\leq i\leq m\right\}.\]
The main problem is that the indicator function $\chi_{[0,\infty)}$ is not a test function.  Moreover, we also need to solve the above variational problem explicitly. It was calculated in \cite{xue2024sample} by using Fourier analysis.  Based on the contraction principle, we present a different approach to solve the above variational problem, which simplifies the proof in \cite{xue2024sample} significantly. 

\subsubsection{The lower bound} We need to prove that for any open set $\mathcal{O} \subset \R^m$,
\begin{equation*}
	\liminf_{n \rightarrow \infty} \frac{n}{a_n^2} \log \mathbb{P}_{\nu_{\rho}} \Big( \{\bar{J}^n_{-1,0} (t_i)/a_n, 1 \leq i \leq m\} \in \mathcal{O} \Big) \geq - \inf_{\mathbf{r} \in \mathcal{O}} \frac{1}{2} \mathbf{r}^T A^{-1} \mathbf{r}.
\end{equation*}
Using the relation between the current and the fluctuation field in \eqref{current decom} and the super-exponential estimates in \eqref{super exp 1} and \eqref{exp tight error}, for any $\mathbf{r} \in  \mathcal{O}$ and any  $\varepsilon > 0$ such that $B(\mathbf{r},\varepsilon) \subset \mathcal{O}$,
 \begin{align*}
 	\liminf_{n \rightarrow \infty} \frac{n}{a_n^2} \log &\mathbb{P}_{\nu_{\rho}} \Big( \{\bar{J}^n_{-1,0} (t_i)/a_n, 1 \leq i \leq m\} \in \mathcal{O} \Big) \\
 	&\geq - \liminf_{l\rightarrow+\infty} \inf \{\mathcal{Q}^\beta_{t_m} (\mu): (\langle \mu_{t_i} - \mu_{0}, \tilde{G}_l\rangle, 1 \leq i \leq m ) \in B(\mathbf{r},\varepsilon)\}\\
 	&\geq - \liminf_{l\rightarrow+\infty} \inf \{\mathcal{Q}^\beta_{t_m} (\mu): \langle \mu_{t_i} - \mu_{0}, \tilde{G}_l \rangle=r_i\text{~for all~}1\leq i\leq m \},
 \end{align*}
where $B(\mathbf{r},\varepsilon)$ is the ball of radius $\varepsilon$ centered at the point $\mathbf{r}$, and $\tilde{G}_l$ was introduced before \eqref{exp tight error}. We refer the readers to \cite[Proof of (6.2)]{xue2024sample} for details of the above argument.  

To calculate the above infimum, we have the following result.

\begin{lemma}\label{lemma 5.1}
Let $G\in \Scal(\R)$ . For any $m \geq 1$,  any $0 \leq t_1 < t_2 < \ldots < t_m \leq T$ and any $\mathbf{r}=(r_1,\ldots, r_m)^T\in \mathbb{R}^m$,
\[		\inf\left\{\mathcal{Q}^\beta_{t_m} (\mu):~ \langle \mu_{t_i} - \mu_{0}, G \rangle=r_i\text{~for all~}1\leq i\leq m\right\}
=\sup_{\bm{\xi} \in \R^m}\left\{\bm{\xi}^T \mathbf{r}-\frac{1}{2}\bm{\xi}^T \Sigma \bm{\xi}\right\},\]
	where $\Sigma := \Sigma_{\{t_k\}_{k=1}^m}(G)$ is a $m\times m$ symmetric matrix such that
	\[
	\Sigma (i,j)={\rm Cov}\left(\mathcal{Y}_{t_i}(G)-\mathcal{Y}_{0}(G), \mathcal{Y}_{t_j}(G)-\mathcal{Y}_{0}(G)\right)
	\]
	for any $1\leq i,j\leq m$. In particular, when $\Sigma$ is invertible, 
	\[
\inf\left\{\mathcal{Q}^\beta_{t_m} (\mu):~ \langle \mu_{t_i} - \mu_{0}, G \rangle=r_i\text{~for all~}1\leq i\leq m\right\}=\frac{1}{2} \mathbf{r}^T\Sigma^{-1} \mathbf{r} .
	\]
\end{lemma}

\begin{proof}
	Since $\{\mathcal{Y}_t\}_{t\geq 0}$ is Gaussian,   the vector $\{\frac{\sqrt{n}}{a_n}\mathcal{Y}_{t_i}(G)-\frac{\sqrt{n}}{a_n}\mathcal{Y}_{0}(G):~1\leq i\leq m\}_{n\geq 1}$ satisfies the large deviation principles  with rate function 
	\[\sup_{\bm{\xi}\in \R^m}\left\{\bm{\xi}^T \mathbf{r}-\frac{1}{2} \bm{\xi}^T \Sigma \bm{\xi} \right\},\]
	see \cite{deuschel1989large} for example. Then, by Lemma \ref{lem MPD of the fluctuation limit} and the contraction principle, the first identity in Lemma \ref{lemma 5.1} holds. When $\Sigma$ is invertible, the second identity in Lemma \ref{lemma 5.1} follows directly from 
	Cauchy-Schwarz inequality. 
\end{proof}

By direct calculations, it is straightforward to prove that  for any $s,t > 0$,
\begin{equation}\label{cov converg}
\lim_{l\rightarrow+\infty}{\rm Cov}\left(\mathcal{Y}_t(\tilde{G}_l)-\mathcal{Y}_0(\tilde{G}_l), \mathcal{Y}_s(\tilde{G}_l)-\mathcal{Y}_0(\tilde{G}_l)\right) = a(t,s),
\end{equation}
where $a(\cdot,\cdot)$ was defined in \eqref{a t s}. Then, by Lemma \ref{lemma 5.1}, 
\begin{align*}
	&\liminf_{n \rightarrow \infty} \frac{n}{a_n^2} \log \mathbb{P}_{\nu_{\rho}} \Big( \{\bar{J}^n_{-1,0} (t_i)/a_n, 1 \leq i \leq m\} \in \mathcal{O} \Big) \\
	&\geq - \liminf_{l\rightarrow+\infty}  \frac{1}{2} \mathbf{r}^T \big(\Sigma_{\{t_k\}_{k=1}^m}(\tilde{G}_l)\big)^{-1} \mathbf{r} \\
	&=  - \frac{1}{2} \mathbf{r}^T A^{-1} \mathbf{r},
\end{align*}
where the matrix $A$ was defined in Lemma \ref{lem: finiteD mdp current}.  We conclude the proof of the lower bound by optimizing over $\mathbf{r} \in \mathcal{O}$. 

\subsubsection{The upper bound} By exponential tightness of the current, we only need to prove the upper bound for any compact set $\mathcal{K} \subset \R^m$. Following \cite[Proof of (6.1)]{xue2024sample}  line by line, one can show that, for any $\varepsilon > 0$,
\begin{multline*}
	\limsup_{n \rightarrow \infty} \frac{n}{a_n^2} \log \mathbb{P}_{\nu_{\rho}} \Big( \{\bar{J}^n_{-1,0} (t_i)/a_n, 1 \leq i \leq m\} \in \mathcal{K} \Big)\\
	\leq - \inf_{\mathbf{r} \in \mathcal{K}} \inf_{\mu}  \left\{\mathcal{Q}^\beta_{t_m} (\mu):~\int_0^{+\infty}[\mu(t_i,u)-\mu(0, u)]du=r_i\text{~for all~}1\leq i\leq m\right\}.
\end{multline*}

The following result bounds from below the first infimum in the last inequality, from which the upper bound follows immediately.

\begin{lemma}\label{lemma contraction principle of current}
For any $\beta\geq 0$ and any $\mathbf{r}=(r_1,\ldots, r_m)^T\in \mathbb{R}^m$, 
\[
\inf\left\{\mathcal{Q}^\beta_{t_m} (\mu):~\int_0^{+\infty} [\mu(t_i,u)-\mu(0, u)] du=r_i\text{~for all~}1\leq i\leq m\right\}\geq \frac{1}{2} \mathbf{r}^T A^{-1} \mathbf{r},
\]
where the matrix $A$ was defined in Lemma \ref{lem: finiteD mdp current}.
\end{lemma}

In order to prove Lemma \ref{lemma contraction principle of current}, we first calculate the macroscopic current across the origin.

\begin{lemma}\label{lemma 5.3}
	Assume $\mu\in D ([0, T], \Scal^\prime(\R))$ satisfies that $\mathcal{Q}_{T}^\beta (\mu)<+\infty$.   Let  $G = G(\mu)$ and $\varphi = \varphi (\mu)$ be identified  in Lemma \ref{lemma finite rate}.
	\begin{enumerate}[(i)]
		\item For $\beta<1$,
		\[
		\int_0^{+\infty}[\mu(t,u)-\varphi (u)]du=\int_{-\alpha(1-2\rho)t}^0 \varphi (u) du, \quad 0\leq t\leq T.
		\]
		\item For $\beta>1$,
		\begin{align*}
			\int_0^{+\infty}&[\mu(t,u)-\varphi (u)]du\\
			&=\int_{\R}\varphi (u) V_t(u)du-\int_0^t\left(\int_{\R}\partial_{uu}^2G (s,u)\mathbb{P}(B_{t-s}+u\geq 0)du\right)ds, \quad 0 \leq t \leq T,
		\end{align*}
		where $\{B_t\}_{t\geq 0}$ is the one dimensional standard Brownian motion starting from the origin and 
		\[
		V_t(u)=\mathbb{P}(B_t+u\geq 0)-\chi_{\{u\geq 0\}}=
		\begin{cases}
			\mathbb{P}(B_t\geq |u|) & \text{~if~}u<0,\\
			-\mathbb{P}(B_t\geq |u|) & \text{~if~}u\geq 0. 
		\end{cases}
		\]
		\item For $\beta=1$,
		\begin{align*}
			&\int_0^{+\infty} [\mu(t,u)-\varphi (u)] du=\int_{\R}\varphi (u)R_t(u)du\\
			&-\int_0^t\left(\int_{\R}\partial_{uu}^2 G (s,u)\mathbb{P}(B_{t-s}+u+\alpha(t-s)(1-2\rho)\geq 0)du\right)ds, \quad 0 \leq t \leq T,
		\end{align*}
		where 
		\[
		R_t(u)=\mathbb{P}(B_t\geq -u-\alpha(1-2\rho)t)-\chi_{\{u\geq 0\}}=
		\begin{cases}
			\mathbb{P}(B_t\geq -u-\alpha(1-2\rho)t) & \text{~if~}u<0,\\
			-\mathbb{P}(B_t\geq u+\alpha(1-2\rho)t) & \text{~if~}u\geq 0. 
		\end{cases}
		\] 
	\end{enumerate}
\end{lemma}

\begin{remark}
	Note that the integral in (ii) converges since
	\begin{align*}
		\int_{\R}\partial_{uu}^2G (s,u)\mathbb{P}(B_{t-s}+u\geq 0)du&=-\int_{\R}\partial_{u}G (s,u)\partial_u\mathbb{P}(B_{t-s}+u\geq 0)du\\
		&=-\int_{\R}\partial_{u}G (s,u)\frac{1}{\sqrt{2\pi(t-s)}}e^{-\frac{u^2}{2(t-s)}}du.
	\end{align*}
	Similarly, the integral in (iii) also converges. 
\end{remark}


	\begin{proof}
		
		We only deal with the  case $\beta = 1$ and the remaining two cases follow from the same analysis. Since $\mathcal{Q}^\beta_T (\mu) < + \infty$, by Lemma \ref{lemma finite rate},
		\[
		\partial_t\mu(t,u)=\frac{1}{2}\Delta\mu(t, u)-\alpha(1-2\rho)\nabla\mu(t, u)-\Delta G (t,u).
		\]
		Thus, for $0\leq t\leq T$,
		\[
		\mu(t,u)=S_{-t}T_t \varphi (u)-\int_0^tS_{-(t-s)}T_{t-s}\Delta G (s,u)ds,
		\]
		where $\{T_t\}_{t\geq 0}$ is the semigroup associated with the Laplacian $(1/2) \Delta$ and $\{S_s\}_{s \in \R}$ is the translation operator: $S_s h (u) = h (u+\alpha(1-2\rho)s)$. Then, 
		\[
		\int_0^{+\infty} [\mu(t,u)- \varphi (u)]du={\rm \uppercase\expandafter{\romannumeral1}}-{\rm \uppercase\expandafter{\romannumeral2}},
		\]
		where
		\begin{align*}
			{\rm \uppercase\expandafter{\romannumeral1}} &=\int_0^{+\infty}S_{-t}T_t\varphi (u)-\varphi (u)du,\\
			{\rm \uppercase\expandafter{\romannumeral2}} &=\int_0^{+\infty}\left(\int_0^tS_{-(t-s)}T_{t-s}\Delta G(s,u)ds\right)du.
		\end{align*}
		By direct calculations,
		\begin{align}\label{equ appendix C1}
			{\rm \uppercase\expandafter{\romannumeral1}}&=\lim_{M\rightarrow+\infty}\int_0^M\mathbb{E} [\varphi (B_t-\alpha(1-2\rho)t+u)] -\varphi (u)du\notag\\
			&=\lim_{M\rightarrow+\infty}\int_0^M\left(\int_{-\infty}^{+\infty}\frac{\varphi (v)}{\sqrt{2\pi t}}e^{-\frac{(v-u+\alpha(1-2\rho)t)^2}{2t}}dv\right)-\varphi (u)du \notag\\
			&=\lim_{M\rightarrow+\infty}\int_{-\infty}^{+\infty}\varphi (u)\left(\int_0^M\frac{1}{\sqrt{2\pi t}}e^{-\frac{(v-u-\alpha(1-2\rho)t)^2}{2t}}dv-\chi_{\{0\leq u\leq M\}}\right)du \\
			&=\lim_{M\rightarrow+\infty}\int_{-\infty}^{+\infty}\varphi (u) \left(\mathbb{P}(0\leq B_t+u+\alpha(1-2\rho)t\leq M)-\chi_{\{0\leq u\leq M\}}\right)du \notag\\
			&=\int_{-\infty}^{+\infty}\varphi (u)R_t(u)du \notag,
		\end{align}
		and 
		\begin{align}\label{equ appendix C2}
			{\rm \uppercase\expandafter{\romannumeral2}}&=\int_0^{+\infty}\left(\int_0^t\mathbb{E} [\Delta G (s, B_{t-s}+u-\alpha(1-2\rho)(t-s)) ] ds\right)du
			\notag\\
			&=\int_0^{+\infty}\left(\int_0^t\left(\int_{-\infty}^{+\infty}\frac{\Delta G (s,v)}{\sqrt{2\pi(t-s)}}e^{-\frac{(v-u+\alpha(1-2\rho)(t-s))^2}{2(t-s)}}dv\right)ds\right)du \\
			&=\int_0^t\left(\int_{-\infty}^{+\infty}\Delta G (s,u)\left(\int_0^{+\infty}\frac{1}{\sqrt{2\pi(t-s)}}e^{-\frac{(v-u-\alpha(1-2\rho)(t-s))^2}{2(t-s)}}dv\right)du\right)ds\notag\\
			&=\int_0^t\left(\int_{-\infty}^{+\infty}\Delta G (s,u)\mathbb{P}\left(B_{t-s}+u+\alpha(1-2\rho)(t-s)\geq 0\right)du\right)ds. \notag
		\end{align}
		This concludes the proof.
	\end{proof}


Now, we are ready to prove Lemma \ref{lemma contraction principle of current}. 

\begin{proof}[Proof of Lemma \ref{lemma contraction principle of current}]
	We only deal with the  case $\beta = 1$ since the remaining cases follow from the same analysis.   It suffices to show that, if  $\mathcal{Q}^\beta_{t_m}(\mu)<+\infty$ and $\int_0^{+\infty} [\mu(t_i,u)-\mu(0, u) ] du=r_i$ for all $1\leq i\leq m$, then 
	\[
	\mathcal{Q}^\beta_{t_m} (\mu)\geq \frac{1}{2} \mathbf{r}^T A^{-1} \mathbf{r}.
	\]
	
	We first prove the above inequality for $\mu$ such that
	\begin{equation}\label{equ condition 5.1}
		\varphi \in C^\infty_c(\R) \text{~and~} G \in C_c^{1,+\infty}([0, T]\times \R),
	\end{equation}
where $\varphi$ and $G$ were identified in Lemma \ref{lemma finite rate}. Let $\tilde{G}_l\in \Scal(\R)$ be the approximation of $G_l$ introduced before \eqref{exp tight error}. By Lemma \ref{lemma 5.1},
	\[
	\mathcal{Q}^\beta_{t_m} (\mu)\geq \frac{1}{2} (\mathbf{r}^l)^T\left(\Sigma_{\{t_k\}_{k=1}^m}(\tilde{G}_l)\right)^{-1}\mathbf{r}^l,
	\]
	where $\mathbf{r}^l=(r_1^l,\ldots, r_m^l)^T \in \R^m$ with $r_i^l=\int_{\R}(\mu(t_i, u)-\varphi (u))\tilde{G}_l(u)du$ for all $1\leq i\leq m$.
	We claim that
	\begin{equation}\label{equ claim proof of Lemma 5.2}
		\lim_{l\rightarrow+\infty} r_i^l= r_i.
	\end{equation}
Then, by \eqref{cov converg},
	\[
	\mathcal{Q}_{t_m} (\mu)\geq \limsup_{l\rightarrow+\infty} \frac{(\mathbf{r}^l)^T\left(\Sigma_{\{t_k\}_{k=1}^m}(\tilde{G}_l)\right)^{-1}\mathbf{r}^l}{2} = \frac{1}{2} \mathbf{r}^T A^{-1} \mathbf{r}.
	\]

	Now we prove \eqref{equ claim proof of Lemma 5.2}.  
	As in the proof of Lemma \ref{lemma 5.3}, for any $H\in L^2(\R)$, 
	\begin{align}\label{equ 5.3}
		\int_{\R}\left(\mu(t,u)-\varphi (u)\right)H(u)du=&\int_{\R}\varphi (u)\left(\mathbb{E}[H(B_t+u+\alpha(1-2\rho)t)]-H(u)\right)du \\
		&-\int_0^t\left(\int_{\R}\partial_{uu}^2G (s,u)\mathbb{E}[H\left(B_{t-s}+\alpha(1-2\rho)(t-s)+u\right)] du\right)ds\notag. 
	\end{align}
	By Equation \eqref{equ 5.3}, Cauchy-Schwarz inequality and the fact that $G \in C_c^{1,+\infty}([0, T]\times \R)$, 
	\[
	\lim_{l\rightarrow+\infty}\int_{\R}(\mu(t, u)- \varphi (u))(\hat{G}_l(u)-G_l(u))du=0. 
	\]
	Hence, to check Equation \eqref{equ claim proof of Lemma 5.2}, we only need to show that
	\begin{equation}\label{equ 5.4}
		\lim_{l\rightarrow+\infty}\int_{\R}(\mu(t, u)-\varphi (u)) G_l (u)du=\int_0^{+\infty}\mu(t, u)-\varphi (u) du.
	\end{equation}
	By Lemma \ref{lemma 5.3} and Equation \eqref{equ 5.3},
	\begin{align}\label{equ 5.5}
		&\int_{\R}(\mu(t, u)-\mu(0,u))G_l(u)du-\int_0^{+\infty}\mu(t, u)-\mu(0,u)du \notag\\
		&=\int_{\R} \varphi (u) F_l(u)du-\int_0^t\left(\int_{\R}\partial_{uu}^2 G (s,u)Q_l(s,u)du\right)ds,
	\end{align}
	where
	\begin{align*}
		F_l(u) &=\mathbb{E}[G_l(B_t+u+\alpha(1-2\rho)t)]-G_l(u)-R_t(u),\\
		Q_l(s,u)&=\mathbb{E}[G_l\left(B_{t-s}+\alpha(1-2\rho)(t-s)+u\right)]-\mathbb{P}(B_{t-s}+u+\alpha(t-s)(1-2\rho)\geq 0).
	\end{align*}
	According to the definition of $G_l$, we have
	\begin{align*}
		|Q_l(s,u)|&\leq \frac{\mathbb{E}|B_{t-s}+u+\alpha(1-2\rho)(t-s)|}{l}+\mathbb{P}\left(B_{t-s}+u+\alpha(t-s)(1-2\rho)\geq l\right).  
	\end{align*}
	Therefore, $\lim_{l\rightarrow+\infty}Q_l (s,u)=0$ uniformly on any compact set in $ [0, t]\times \R$. Similarly, $\lim_{l\rightarrow+\infty}F_l=0$ uniformly on any compact set in $\R$. Then, Equation \eqref{equ 5.4} holds according to Equation \eqref{equ 5.5} and the fact that $\varphi$ and $G$ have compact support. 
	
	For $\mu$ not satisfying \eqref{equ condition 5.1}, the idea is to approximate $\varphi (\mu)$ and $G (\mu)$ by $\varphi^\varepsilon \in C^\infty_c (\R)$ and $G^{\varepsilon} \in C_c^{1,\infty} ([0,T] \times \R)$ respectively.  Precisely speaking, let  $\varphi^\varepsilon  \rightarrow \varphi$ in $L^2(\R)$ and $G^{\varepsilon}  \rightarrow G$ in $\mathbb{H}$ as $\varepsilon \rightarrow 0$.  Let $\mu^\varepsilon \in D ([0, t_m], \Scal^\prime(\R))$ be given in Lemma \ref{lemma finite rate} associated with $\varphi^\varepsilon $ and $G^{\varepsilon}$. 
	Then, by Lemma \ref{lemma finite rate},
	\begin{align}\label{equ limit in proof of Lemma 5.2}
		\lim_{\varepsilon\rightarrow 0}\mathcal{Q}^\beta_{t_m}(\mu^\varepsilon)&=  \lim_{\varepsilon\rightarrow 0} \frac{1}{2\chi (\rho)} \left(\|\varphi^\varepsilon\|_{L^2(\R)}^2+ [G^{\varepsilon},G^{\varepsilon}]\right) \notag\\
		&=\frac{1}{2\chi (\rho)} \left(\|\varphi\|_{L^2(\R)}^2+ [G,G]\right)=\mathcal{Q}_{t_m}(\mu)
	\end{align}
	Since $\mu^\varepsilon$ satisfies \eqref{equ condition 5.1},  we have shown that
	\[
	\mathcal{Q}^\beta_{t_m} (\mu^\varepsilon)\geq \frac{1}{2} (\mathbf{r}^\varepsilon)^T A^{-1} \mathbf{r}^\varepsilon,
	\]
	where $\mathbf{r}^\varepsilon=(r^\varepsilon_1,\ldots, r^\varepsilon_m)^T$ with $r^\varepsilon_i=\int_0^{+\infty}\mu^\varepsilon(t_i,u)-\varphi^\varepsilon (u) du$ for all $1\leq i\leq m$.  Thus,
	\[\mathcal{Q}_{t_m}^\beta (\mu) \geq    \limsup_{\varepsilon\rightarrow 0} \frac{1}{2} (\mathbf{r}^\varepsilon)^T A^{-1} \mathbf{r}^\varepsilon.\]
By Lemma \ref{lemma 5.3},  $\lim_{\varepsilon \rightarrow 0} \mathbf{r}^\varepsilon = \mathbf{r}$, thus concluding the proof. 
\end{proof}



\subsection{MDP for the tagged particle}\label{subsec:tagged} As for the current, we only need to prove exponential tightness and finite dimensional MDP for the tagged particle. Since the orderings of particles are preserved by the dynamics, the tagged particle and the current are related as follows: for $x > 0$, 
\begin{align}\label{tagged current relation}
 \{X^n (t) > x\} = \Big\{J^n_{-1,0} (t) \geq \sum_{y=0}^{x} \eta_y (t) \Big\},  \quad \{X^n (t) < -x\} = \Big\{J^n_{-1,0} (t) <  \sum_{y=-x}^{-1} \eta_{y} (t) \Big\}.
\end{align}
Equivalently, 
\begin{align}
J^n_{-1,0} (t) &= \sum_{x=0}^{X^n(t)-1} \eta_x (t) = \sum_{x=0}^{X^n(t)-1} (\eta_x (t) -\rho) + \rho X^n(t) \quad \text{if}\; J^n_{-1,0} (t) \geq 0;\label{tagged current relation 1}\\
J^n_{-1,0} (t) &= -\sum_{x=X^n(t)}^{-1} \eta_x (t) = - \sum_{x=X^n(t)}^{-1} (\eta_x (t) -\rho) + \rho X^n(t) \quad \text{if}\; J^n_{-1,0} (t) < 0.\label{tagged current relation 2}
\end{align}
Following the proof of \cite[Lemma 5.3]{xue2024sample} and using \eqref{tagged current relation}, one can show that the tagged particle process is exponentially tight. Following the proof of \cite[Lemma 6.2]{xue2024sample} and using \eqref{tagged current relation 1}, \eqref{tagged current relation 2}, one can prove finite dimensional MDP for the tagged particle process. Since the proof is exactly the same, we do not repeat it here. Intuitively, the first terms on the right hand sides of \eqref{tagged current relation 1} and \eqref{tagged current relation 2}, divided by $a_n$, are superexponentially small.  Then, it is straightforward to see that $\mathcal{X}^\beta = \rho^2 \mathcal{J}^\beta$ since the rate functions have quadratic forms. 






\appendix

%\section{Proof of Theorem \ref{thm clt filed}}\label{app: pf clt fluc field}
%
%In this section, we outline the proof of Theorem \ref{thm clt filed}. By Dynkin's martingale formula, for any $H \in \Scal (\R)$, 
%\begin{align*}
%    \Mcal^n_t (H) &= \Ycal^n_t (H) - \Ycal^n_0 (H) - \int_0^t L_n \Ycal^n_s (H) ds,\\
%    \big[\Mcal^n_t (H)\big]^2 &- \int_0^t \big\{L_n \Ycal^n_s (H)^2 - 2 \Ycal^n_s (H) L_n \Ycal^n_s (H)\big\} ds
%\end{align*}
%are both martingales.  In the rest of calculations, we write $\eta (s)$ simply as $\eta$ since there is no confusion. Direct calculations show that 
%\begin{align*}
%    L_n \Ycal^n_s (H) = \frac{n^{\gamma-2}}{2} \mathcal{Y}^n_t (\Delta_n H) + \frac{\alpha n^{\gamma-\beta-1}}{\sqrt{n}} \sum_{x \in \Z} \nabla_n H (\tfrac{x}{n}) \big[ j_{x,x+1} (\eta) - \rho(1-\rho)\big].
%\end{align*}
%Above, $\Delta_n$ and $\nabla_n$ are the discrete Laplacians and gradient respectively, 
%\[\Delta_nH(\tfrac{x}{n}) = n^2 \big[ H(\tfrac{x+1}{n}) + H(\tfrac{x-1}{n}) - 2 H(\tfrac{x}{n})   
% \big],  \quad \nabla_nH(\tfrac{x}{n}) = n \big[ H(\tfrac{x+1}{n})  -  H(\tfrac{x}{n})   
% \big],\]
%and $j_{x,x+1} (\eta)= \eta_x (1-\eta_{x+1})$ is the instantaneous current from $x$ to $x+1$ corresponding to the generator $L_a$. Since
%\[j_{x,x+1} (\eta)- \rho(1-\rho) = (1-2\rho) \Bar{\eta}_x + \rho (\Bar{\eta}_x - \Bar{\eta}_{x+1}) -\Bar{\eta}_x \Bar{\eta}_{x+1},\]
%using the summation by parts formula, we have
%\begin{align}
%       L_n \Ycal^n_s (H) =& \frac{n^{\gamma-2}}{2} \mathcal{Y}^n_t (\Delta_n H)   + \alpha (1-2\rho) n^{\gamma-\beta-1}  \Ycal^n_s (\nabla_n H) \nonumber\\
%       &+ \alpha \rho n^{\gamma-\beta-2} \Ycal^n_s (\Delta_n H) - \frac{\alpha n^{\gamma-\beta-1}}{\sqrt{n}} \sum_{x \in \Z} \nabla_n H (\tfrac{x}{n}) \Bar{\eta}_x \Bar{\eta}_{x+1}.\label{ln ys}
%\end{align}
%Similarly, 
%\begin{align}
%    L_n \Ycal^n_s (H)^2 - 2 \Ycal^n_s (H) L_n \Ycal^n_s (H) =& \frac{n^{\gamma-3}}{2} \sum_{x \in \Z} (\eta_x-\eta_{x+1})^2 \big[ \nabla_n H (\tfrac{x}{n})\big]^2 \notag\\
%    &+ \alpha n^{\gamma-\beta-3} \sum_{x \in \Z} \eta_x (1-\eta_{x+1})\big[ \nabla_n H (\tfrac{x}{n})\big]^2. \label{ln ys2}
%\end{align}
%
%By standard arguments, one can show that the sequences $\{\Mcal^n_t, 0 \leq t \leq T\}_{n \geq 1}$ and $\{\Ycal^n_t, 0 \leq t \leq T\}_{n \geq 1}$ are tight in the space $D([0,T], \mathcal{S} (\R))$. Let $\{\Mcal_t, 0 \leq t \leq T\}$ and $\{\Ycal_t, 0 \leq t \leq T\}$ be any limit point along some subsequence. Then, we only need to characterize those limit points. 
%
%Since $\gamma - \beta - 1 \leq 0$, the third term in \eqref{ln ys} and the second one in \eqref{ln ys2} vanish as $n \rightarrow \infty$. For the last term on the right-hand side of \eqref{ln ys}, we have the following result, which is called Boltzmann-Gibbs principle in the literature.  The proof of \cite[Theorem 2.6]{gonccalves2008central} can be adapted here directly and thus is omitted. 
%
%\begin{proposition} For any $H \in \Scal (\R)$,
%   \[\lim_{n \rightarrow \infty} \E_{\nu_\rho} \Big[ \sup_{0 \leq t \leq T} \Big( \int_0^t  \frac{ n^{\gamma-\beta-1}}{\sqrt{n}} \sum_{x \in \Z} \Bar{\eta}_x (s)\Bar{\eta}_{x+1} (s) H(\tfrac{x}{n})ds \Big)^2\Big] = 0.\]  
%\end{proposition}
%
%
%
%For the remaining terms, we need to discuss different regimes of $\beta$.
%
%\medspace
%
%{\bf The case $\beta > 1$.} Since $\gamma=2$, we can replace the first term on the right-hand side of \eqref{ln ys} by $\Ycal_t (\Delta H)/2$, and the second one vanishes as $n \rightarrow \infty$. By Cauchy-Schwarz inequality, we can replace the first term on the right-hand side of \eqref{ln ys2} by $\chi (\rho) \|\nabla H\|_{L^2 (\R)}^2$. Thus, the  limit points $\{\Mcal_t, 0 \leq t \leq T\}$ and $\{\Ycal_t, 0 \leq t \leq T\}$ satisfy
%\begin{align*}
%    \Mcal_t (H) &= \Ycal_t (H) - \Ycal_0 (H) - \frac{1}{2}\int_0^t \Ycal_s (\Delta H) ds,\\
%    \big[\Mcal_t (H)\big]^2 &- t \chi(\rho) \|\nabla H\|_{L^2 (\R)}^2
%\end{align*}
%are both martingales.
%
%\medspace
%
%{\bf The case $\beta = 1$.} Compared with the supercritical case, the second term on the right-hand side of \eqref{ln ys} does not vanish and can be replaced by $\alpha (1-2 \rho) \Ycal_s (\nabla H)$. Thus, 
%\begin{align*}
%    \Mcal_t (H) &= \Ycal_t (H) - \Ycal_0 (H) - \int_0^t \Big\{ \frac{1}{2} \Ycal_s (\Delta H) + \alpha (1-2 \rho) \Ycal_s (\nabla H) \Big\} ds,\\
%    \big[\Mcal_t (H)\big]^2 &- t \chi(\rho) \|\nabla H\|_{L^2 (\R)}^2
%\end{align*}
%are both martingales.
%
%\medspace
%
%{\bf The case $\beta < 1$.} Since $\gamma=1+\beta$, the first term on the right-hand side of \eqref{ln ys} vanishes and the second one can be replaced by $\alpha (1-2 \rho) \Ycal_s (\nabla H)$. The first term on the right-hand side of \eqref{ln ys2} vanishes, and so does the martingale term. Therefore, 
%\begin{align*}
% 0 = \Ycal_t (H) - \Ycal_0 (H) - \alpha (1-2 \rho) \int_0^t   \Ycal_s (\nabla H)  ds.
%\end{align*}
%
%\medspace
%
%Finally, we conclude the proof of Theorem \ref{thm clt filed} by the uniqueness of the solution to the above martingale problems.

%\section{Proof of Equations \eqref{equ covariance sub}-\eqref{equ covariance cri}} \label{Appendix A.1}
%In this subsection, we prove Equations \eqref{equ covariance sub}-\eqref{equ covariance cri}. We only check the  case $\beta > 1$ and other two cases follow from the same analysis.
%By solving the equation
%	\[
%	\partial_t \Ycal_t = \frac{1}{2} \Delta \Ycal_t + \sqrt{\chi (\rho)} \nabla d \mathcal{W}_t, 
%	\]
%	we have
%	\begin{equation*}
%		\mathcal{Y}_t=T_t\mathcal{Y}_0+\sqrt{\chi(\rho)}\int_0^tT_{t-\theta}\nabla d\mathcal{W}_\theta
%	\end{equation*}
%	for any $t\geq 0$. Hence, for $t \geq s>0$ and $H,G\in \Scal (\R)$,
%	\begin{equation}\label{equation A.1.1}
%		{\rm Cov}(\mathcal{Y}_t(H), \mathcal{Y}_s(G))
%		=\chi(\rho)\langle T_tH, T_sG\rangle+\chi(\rho)\int_0^s \langle \nabla T_{t-\theta}H, \nabla T_{s-\theta}G\rangle d\theta.
%	\end{equation}
%	For $0\leq \theta\leq s$, let
%	\[
%	g(\theta)=\langle T_{t-\theta}H, T_{s-\theta}G\rangle,
%	\]
%	then
%	\begin{align*}
%		\frac{d}{d\theta}g(\theta)&=-\frac{1}{2}\langle \Delta T_{t-\theta}H, T_{s-\theta}G\rangle-\frac{1}{2}\langle  T_{t-\theta}H, \Delta T_{s-\theta}G\rangle\\
%		&=\frac{1}{2}\langle \nabla T_{t-\theta}H, \nabla T_{s-\theta}G\rangle+\frac{1}{2}\langle \nabla T_{t-\theta}H, \nabla T_{s-\theta}G\rangle\\
%		&=\langle \nabla T_{t-\theta}H, \nabla T_{s-\theta}G\rangle.
%	\end{align*}
%	As a result,
%	\begin{align*}
%		\int_0^s \langle \nabla T_{t-\theta}H, \nabla T_{s-\theta}G\rangle d\theta = g(s)-g(0) = \langle T_{t-s}H, G \rangle-\langle T_tH, T_sG\rangle,
%	\end{align*}
%	concluding the proof of  Equation \eqref{equ covariance sup}.



%\section{Proof of Lemma \ref{lem MPD of the scaled density field}}
%
%
%\subsection{The upper bound} For $H \in \mathcal{C}^{1,\infty}_c ([0,T] \times \R)$,   by Feynman-Kac formula \cite[Appendix 1.7]{klscaling},
%\begin{align*}
%	\mathfrak{M}^n_t (H) = \exp \Big\{   \frac{a_n^2}{n}  \langle \mu^n_t ,H_t\rangle -  \frac{a_n^2}{n}  \langle \mu^n_0 ,H_0\rangle - \int_0^t e^{-\frac{a_n^2}{n}  \langle\mu^n_s,H_s\rangle}(\partial_{s} +L_n) e^{\frac{a_n^2}{n}  \langle\mu^n_s ,H_s\rangle} ds \Big\}
%\end{align*}
%is a martingale.  
%
%We first calculate the time integral inside the above exponential. For the symmetric part, by Taylor's expansion to the second order,
%\begin{equation}\label{cal 4}
%	\begin{aligned}
%		&e^{-\frac{a_n^2}{n}  \langle \mu^n_s,H_s\rangle} n^\gamma L_s e^{\frac{a_n^2}{n}  \langle \mu^n_s ,H_s\rangle} 
%		=  \frac{n^\gamma}{2} \sum_{x \in \Z}  \Big[e^{\frac{a_n}{n} (\eta_x - \eta_{x+1} ) (H_s(\frac{x+1}{n}) - H_s (\frac{x}{n}))} - 1\Big] \\
%		& = \frac{a_n^2}{n} \Big\{  \frac{n^{\gamma-2}}{2} \langle \mu^n_s,\Delta_n H_s\rangle + \frac{n^{\gamma-2}}{4n} \sum_{x \in \Z} (\eta_x - \eta_{x+1})^2 \big(\nabla_{n} H_s(\tfrac{x}{n})\big)^2 + o_n (1) \Big\},
%	\end{aligned}
%\end{equation}
%where $o_n (1)$ converges to zero as $n \rightarrow \infty$.  Here, $\nabla_{n}$ and $\Delta_n$ are the discrete gradient and Laplacian respectively,
%\begin{align*}
%	\nabla_{n} H_s(\tfrac{x}{n}) &= n \big[H_s(\tfrac{x+1}{n}) - H_s(\tfrac{x}{n})\big], \\
%	\Delta_{n} H_s(\tfrac{x}{n}) &= n^2 \big[H_s(\tfrac{x+1}{n}) + H_s(\tfrac{x-1}{n}) - 2H_s(\tfrac{x}{n})\big].
%\end{align*}
%The asymmetric part can be calculated similarly: by Taylor's expansion to the first order,
%\begin{equation}\label{cal 5}
%	\begin{aligned}
%		&e^{-\frac{a_n^2}{n}  \langle \mu^n_s,H_s\rangle} \alpha n^{\gamma-\beta} L_a e^{\frac{a_n^2}{n}  \langle \mu^n_s ,H_s\rangle}   = \alpha n^{\gamma-\beta}  \sum_{x \in \Z}    \eta_x (1-\eta_{x+1}) \Big[e^{\frac{a_n}{n}  (H_s(\frac{x+1}{n}) - H_s (\frac{x}{n}))} - 1\Big]\\
%		&= \frac{a_n^2}{n} \Big\{ \frac{\alpha n^{\gamma-\beta-1}}{a_n}  \sum_{x \in \Z}  [\eta_x (1-\eta_{x+1}) - \chi(\rho) ]\nabla_{n} H_s(\tfrac{x}{n}) + o_n (1) \Big\}\\
%		&= \frac{a_n^2}{n} \Big\{\alpha (1-2\rho)n^{\gamma-\beta-1}\langle \mu^n_s, \nabla_n H_s \rangle -  \frac{\alpha n^{\gamma-\beta-1}}{a_n} \sum_{x \in \Z} \bar{\eta}_x \bar{\eta}_{x+1} \nabla_{n} H_s(\tfrac{x}{n}) + o_n (1) \Big\}.
%	\end{aligned}
%\end{equation}
%We introduced the constant $\chi(\rho)$ into the penultimate line since $\sum_{x \in \Z} \nabla_{n} H_s(\tfrac{x}{n})  = 0$.  In the last identity, we used the fact
%\[\eta_x (1-\eta_{x+1}) - \chi (\rho) = (1-2\rho) \bar{\eta}_x + \rho (\bar{\eta}_x - \bar{\eta}_{x+1}) - \bar{\eta}_x \bar{\eta}_{x+1},\]
%and note that, since $\gamma \leq \beta+1$,
%\[\frac{\alpha n^{\gamma-\beta-1}}{a_n}  \sum_{x \in \Z} (\bar{\eta}_x - \bar{\eta}_{x+1})  \nabla_{n} H_s(\tfrac{x}{n})  = \frac{\alpha n^{\gamma-\beta-2}}{a_n}  \sum_{x \in \Z} \bar{\eta}_x  \Delta_{n} H_s(\tfrac{x}{n}) = o_n (1).\]
%Finally, for the time derivative, 
%\begin{align*}
%	&e^{-\frac{a_n^2}{n}  \langle \mu^n_s,H_s \rangle} \partial_{s} e^{\frac{a_n^2}{n}  \langle \mu^n_s ,H_s \rangle } = \langle \mu^n_s, \partial_s H_s \rangle.
%\end{align*}
%To sum up, we have shown that
%\begin{equation}\label{exp martingale}
%	\begin{aligned}
%		\mathfrak{M}^n_t (H) &= \exp \Big\{   \frac{a_n^2}{n}  \Big[  \langle \mu^n_t ,H_t\rangle -   \langle \mu^n_0 ,H_0\rangle \\
%		&- \int_0^t \frac{n^{\gamma-2}}{2} \langle \mu^n_s,\Delta_n H_s\rangle + \langle \mu^n_s, \partial_s H_s \rangle + \alpha (1-2\rho)n^{\gamma-\beta-1}\langle \mu^n_s, \nabla_n H_s \rangle ds \\
%		&- \int_0^t  \frac{n^{\gamma-2}}{4n} \sum_{x \in \Z} (\eta_x - \eta_{x+1})^2 \big(\nabla_{n} H_s(\tfrac{x}{n})\big)^2 ds \\
%		&+ \int_0^t  \frac{\alpha n^{\gamma-\beta-1}}{a_n}  \sum_{x \in \Z} \bar{\eta}_x \bar{\eta}_{x+1} \nabla_{n} H_s(\tfrac{x}{n}) ds  + t o_n (1)  \Big] \Big\}.
%	\end{aligned}
%\end{equation}
%
%The next step is to write the term inside the exponential in \eqref{exp martingale} as a functional of the fluctuation field $\mu^n_t$. By \cite[Lemma 3.1]{bibid}, we could replace the term $\int_0^t \frac{n^{\gamma-2}}{2} \langle \mu^n_s,\Delta_n H_s\rangle ds$ in \eqref{exp martingale} with $\int_0^t \frac{n^{\gamma-2}}{2} \langle \mu^n_s,\Delta H_s\rangle ds$,  and replace the third line in \eqref{exp martingale} with $\frac{n^{\gamma-2} \chi (\rho)}{2} \| \nabla H\|_{L^2 ([0,t] \times \R)}^2$.  For the last line in \eqref{exp martingale}, we have the following lemma. 
%
%\begin{lemma}\label{lem: quadratic term vanish}
%If aaa, then for any $\delta > 0$,
%\begin{equation}
%	\limsup_{n \rightarrow \infty} \frac{n}{a_n^2} \log \mathbb{P}_{\nu_\rho} \Big( \sup_{0 \leq t \leq T} \Big|  \int_0^t   \frac{\alpha n^{\gamma-\beta-1}}{a_n}  \sum_{x \in \Z} \bar{\eta}_x \bar{\eta}_{x+1} \nabla_{n} H_s(\tfrac{x}{n}) ds \Big| > \delta  \Big) = - \infty.
%\end{equation}
%\end{lemma}
%
%\begin{proof}
%	The proof is similar to \cite[Proposition 3.2]{bibid}, and thus we only sketch it.  For $\ell > 0$, let $p_\ell$ be the uniform measure on $\{0,1,\ldots,\ell-1\}$, and let $q_\ell = p_\ell \ast p_\ell$ be the convolution of $p_\ell$ with itself.  For a configuration $\eta$, define
%	\[\Bar{\eta}^\ell_x = \sum_{x \in \Z} \Bar{\eta}_{x+z}q_\ell (z), \quad x \in \Z.\]
%	In the proof, we shall take 
%	\[\ell = \ell(n) = \Big(\frac{n^{2\beta+4}}{n^\gamma a_n^2}\Big)^{1/3}.\]
%	To prove the lemma, we only need to show that, for any $\delta > 0$,
%	\begin{align}
%		\limsup_{n \rightarrow \infty} \frac{n}{a_n^2} \log  \mathbb{P}_{\nu_\rho} \Big( \sup_{0 \leq t \leq T} \Big|  \int_0^t   \frac{\alpha n^{\gamma-\beta-1}}{a_n}  \sum_{x \in \Z} \big( \bar{\eta}_{x+1} - \Bar{\eta}_{x+1}^\ell \big)  \bar{\eta}_{x} \nabla_{n} H_s(\tfrac{x}{n}) ds \Big| > \delta  \Big) = - \infty,\label{quadratic term 1}\\
%		\limsup_{n \rightarrow \infty} \frac{n}{a_n^2}  \log \mathbb{P}_{\nu_\rho} \Big( \sup_{0 \leq t \leq T} \Big|  \int_0^t   \frac{\alpha n^{\gamma-\beta-1}}{a_n}  \sum_{x \in \Z} \bar{\eta}^\ell_{x+1} \Bar{\eta}_{x} \nabla_{n} H_s(\tfrac{x}{n}) ds \Big| > \delta  \Big) = - \infty.\label{quadratic term 2}
%	\end{align}
%
%We first prove \eqref{quadratic term 1}. By Markov's exponential inequality, it suffices to prove that, for any $A >0$,
%\begin{equation}
%\limsup_{n \rightarrow \infty} \frac{n}{a_n^2}  \log \mathbb{E}_{\nu_\rho} \Big[ \exp \Big\{ \sup_{0 \leq t \leq T} \Big|  \int_0^t   \alpha A a_n n^{\gamma-\beta-2}  \sum_{x \in \Z} \big( \bar{\eta}_{x+1} - \Bar{\eta}_{x+1}^\ell \big)  \bar{\eta}_{x} \nabla_{n} H_s(\tfrac{x}{n}) ds \Big|  \Big\} \Big] = 0.
%\end{equation}
%To remove the supremum over time inside the above exponential, define $t_i = iT/n^2$ for $i = 0,\ldots,n^2$.  Then, using the basic inequality
%\[\log E [\exp \{\max_{1 \leq i \leq K} X_i\}] \leq \log K + \max_{1 \leq i \leq K} \log E[\exp \{X_i\}],\]
%and since $a_n \gg \sqrt{n \log n}$, we only need to show that
%\begin{equation}\label{quadratic term 3}
%	\limsup_{n \rightarrow \infty} \sup_{0 \leq t \leq T} \frac{n}{a_n^2}  \log \mathbb{E}_{\nu_\rho} \Big[ \exp \Big\{  \Big|  \int_0^t   \alpha A a_n n^{\gamma-\beta-2}  \sum_{x \in \Z} \big( \bar{\eta}_{x+1} - \Bar{\eta}_{x+1}^\ell \big)  \bar{\eta}_{x} \nabla_{n} H_s(\tfrac{x}{n}) ds \Big|  \Big\} \Big] = 0.
%\end{equation}
%Note that the absolute value inside can be removed similarly. By Feynman-Kac formula, the last expression is bounded by
%\begin{align*}
%	\frac{\alpha A n^{\gamma-\beta-1}}{a_n} \int_0^t ds \sup_{f: \nu_{\rho}-\text{density}} \Big\{ E_{\nu_\rho} \Big[ \sum_{x \in \Z} \big( \bar{\eta}_{x+1} - \Bar{\eta}_{x+1}^\ell \big)  \bar{\eta}_{x} \nabla_{n} H_s(\tfrac{x}{n})  f\Big] - \frac{n^{\beta+2}}{\alpha A a_n} D (\sqrt{f})\Big\}.
%\end{align*}
%A standard argument yields that, for any $B > 0$,
%\begin{align*}
%	E_{\nu_\rho} \Big[ \sum_{x \in \Z} \big( \bar{\eta}_{x+1} - \Bar{\eta}_{x+1}^\ell \big)  \bar{\eta}_{x} \nabla_{n} H_s(\tfrac{x}{n})  f\Big]   \leq C(H) B^{-1} n \ell + B \ell D(\sqrt{f}).
%\end{align*}
%Taking $B = \frac{n^{\beta+2}}{\alpha A a_n \ell}$, we bound the expression in \eqref{quadratic term 3} by 
%\[\frac{C(H,\alpha,A) t \ell^2 n^\gamma}{n^{2\beta+2}}.\]
%
%For \eqref{quadratic term 2}, we only need to prove that for any $A > 0$,
%\begin{equation}
%\limsup_{n \rightarrow \infty} \frac{n}{a_n^2}  \log \mathbb{E}_{\nu_\rho} \Big[ \exp \Big\{  \int_0^T  \Big|   \alpha A n^{\gamma-\beta-2} a_n  \sum_{x \in \Z} \bar{\eta}^\ell_{x+1} \Bar{\eta}_{x} \nabla_{n} H_s(\tfrac{x}{n}) \Big| ds   \Big\}  \Big] = 0.
%\end{equation}
%As before, we can remove the absolute value inside the above exponential.  Since the process is stationary, by Jensen's inequality,  the last expression is bounded by
%\begin{equation}\label{quadratic term 4}
%\frac{n}{a_n^2}  \log \Big( \frac{1}{T}\int_0^T ds \mathbb{E}_{\nu_\rho} \Big[ \exp \Big\{    \alpha A T  n^{\gamma-\beta-2} a_n  \sum_{x \in \Z} \bar{\eta}^\ell_{x+1} \Bar{\eta}_{x} \nabla_{n} H_s(\tfrac{x}{n})    \Big\}  \Big] \Big). 
%\end{equation}
%By change of variables, we can rewrite
%\begin{align*}
%	\sum_{x \in \Z} \bar{\eta}^\ell_{x+1} \Bar{\eta}_{x} \nabla_{n} H_s(\tfrac{x}{n}) = \sum_{y \in \Z} \big(\sum_{x \in \Z} \bar{\eta}_x p_\ell (y-x-1) \nabla_{n} H_s(\tfrac{x}{n}) \big) \big(\sum_{z \in \Z} \bar{\eta}_z p_\ell (z-y)\big). 
%\end{align*}
%By H{\"o}lder's inequality and \cite[Lemma F.8]{jara2018non}, since $a_n \ll n^{2+\beta-\gamma}$, we bound the expectation in \eqref{quadratic term 4} by
%\begin{align*}
%&\prod_{y \in \Z}\mathbb{E}_{\nu_\rho} \Big[ \exp \Big\{   5  \alpha A T  \ell n^{\gamma-\beta-2} a_n  \big(\sum_{x \in \Z} \bar{\eta}_x p_\ell (y-x-1) \nabla_{n} H_s(\tfrac{x}{n}) \big) \big(\sum_{z \in \Z} \bar{\eta}_z p_\ell (z-y)\big) \Big\}  \Big]^{1/(5\ell)}\\
%&\leq C(A,T,H,\alpha)^{n/\ell}.
%\end{align*}
%Therefore,  \eqref{quadratic term 4} is bounded by a constant multiple of 
%\[\frac{n^2}{\ell a_n^2} = \Big( \frac{n^{2+\gamma-2\beta}}{a_n^4} \Big)^{1/3}.\]
%This concludes the proof.
%\end{proof}

%\section{Proof of Lemma \ref{lem MPD of the fluctuation limit}} \label{sec: mdp fluc limit}
%
%\begin{proof}
%	[Outline of the proof of Lemma \ref{lem MPD of the fluctuation limit}] We utilize an exponential martingale strategy similar with that given in the proof of Lemma \ref{lem MPD of the scaled density field}. Here we only give an outline of the proof to avoid too many similar details. We only deal with the critical case, i.e., $\beta=1$, since other two cases follow from the same analysis. For any $H\in C_c^{1,+\infty}([0, T]\times \mathbb{R})$, we claim that
%	\begin{align}\label{equ continuous martingale}
%		\Bigg\{
%		\exp\Big\{&\frac{a_n}{\sqrt{n}}\mathcal{Y}_t(H_t)-\frac{a_n}{\sqrt{n}}\mathcal{Y}_0(H_0)\\
%		&-\frac{a_n}{\sqrt{n}}\int_0^t \mathcal{Y}_s\left((\partial_s+\mathcal{P})H_s\right)ds-\frac{\chi(\rho)a_n^2}{2n}\int_0^t\langle\nabla H_s, \nabla H_s \rangle ds
%		\Big\}
%		\Bigg\}_{0\leq t\leq T} \notag
%	\end{align}
%	is a martingale. We check this claim at the end of this proof. 
%	
%	With martingale \eqref{equ continuous martingale} playing the role as the exponential martingale ..., the MDP upper bound of $\left\{\frac{\sqrt{n}}{a_n}\mathcal{Y}_t:~0\leq t\leq T\right\}_{n\geq 1}$ for compact sets and the lower bound for open sets follow from analyse similar with those given in the proof of Lemma \ref{lem MPD of the scaled density field} and hence we omit the details.  
%	
%	We still need to check the exponential tightness of $\left\{\frac{\sqrt{n}}{a_n}\mathcal{Y}_t:~0\leq t\leq T\right\}_{n\geq 1}$. Let $\pi$ be the distribution of $\mathcal{Y}_0$, i.e., $\nu(H)$ follows the normal distribution with mean $0$ and variance $\chi(\rho)\langle H, H\rangle$ for any $H\in \Scal(\R)$ and the random element $\nu$ in $\Scal^\prime(\R)$ distributed by $\pi$. By Equation \eqref{equ covariance cri}, $\pi$ is an invariant distribution of $\{\mathcal{Y}_t\}_{t\geq 0}$. Then, by applying the Garsia-Rademich-Rumsey lemma, the check of the exponential tightness of $\left\{\frac{\sqrt{n}}{a_n}\mathcal{Y}_t:~0\leq t\leq T\right\}_{n\geq 1}$ starting from invariant distribution $\pi$ follows from the same analysis as that of $\{\mu^n\}_{n\geq 1}$ starting from invariant distribution $\nu_\rho$. Hence, we omit the details. 
%	
%	At last, we check that \eqref{equ continuous martingale} is a martingale. According to It\^{o}'s formula, for any $H\in C_c^{1,+\infty}([0, T]\times \mathbb{R})$, we only need to check that $\{\Lambda_t^H\}_{0\geq t\leq T}$ is a martingale and
%	\[
%	d\langle \Lambda^H\rangle_t=\chi(\rho)\langle \nabla H_t, \nabla H_t \rangle dt,
%	\]
%	where
%	\[
%	\Lambda_t^H=\mathcal{Y}_t(H_t)-\mathcal{Y}_0(H_0)-\int_0^t\mathcal{Y}_s\left((\partial_s+\mathcal{P})H_s\right)ds
%	\]
%	and $\{\langle \Lambda^H \rangle_t\}_{t\geq 0}$ is the quadratic variation process of $\Lambda_t^H$. Let $\mathcal{B}$ be the set of $H \in C_c^{1,+\infty}([0, T]\times \mathbb{R})$ such that $H(t,u)=b_th(u)$ for all $0\leq t\leq T, u\in \mathbb{R}$ and some $b\in C^1[0, T], h\in C_c^\infty(\mathbb{R})$. Since ${\rm span}\mathcal{B}$ is dense in $C_c^{1,+\infty}([0, T]\times \mathbb{R})$, by Equations \eqref{equ covariance sub}-\eqref{equ covariance cri} and the martingale convergence theorem given in \cite{whitt2007proofs}, we only need to deal with $H\in \mathcal{B}$. 
%	
%	For any $H\in \mathcal{B}$, $\mathcal{Y}_t(H_t)=b_t\mathcal{Y}_t(h)$. By It\^{o}'s formula, we have
%	\[
%	d\mathcal{Y}_t(H_t)=\partial_tb_t\mathcal{Y}_t(h)dt+b_td\mathcal{Y}_t(h)=\mathcal{Y}_t(\partial_tH_t)dt+b_td\mathcal{Y}_t(h).
%	\]
%	According to the definition of $\{\mathcal{Y}_t\}_{t\geq 0}$, we have
%	\[
%	d\mathcal{Y}_t(h)=\mathcal{Y}_t(\mathcal{P}h)dt+d\mathcal{M}_t(h),
%	\]
%	where $\{\mathcal{M}_t(h)\}_{t\geq 0}$ is a martingale such that $d\langle \mathcal{M}(h)\rangle_t=\chi(\rho)\langle \nabla h, \nabla h\rangle dt$. Therefore,
%	\begin{align*}
%		d\mathcal{Y}_t(H_t)&=\mathcal{Y}_t(\partial_tH_t)dt+b_t\mathcal{Y}_t(\mathcal{P}h)dt+b_td\mathcal{M}_t(h) \\
%		&=\mathcal{Y}_t\left((\partial_t+\mathcal{P})H_t\right)dt+b_td\mathcal{M}_t(h)
%	\end{align*}
%	and hence $d\Lambda_t^H=b_td\mathcal{M}_t(h)$. Consequently, $\{\Lambda_t^H\}_{0\leq t\leq T}$ is a martingale such that
%	\[
%	d\langle \Lambda^H \rangle_t=\chi(\rho)b_t^2\langle \nabla h, \nabla h\rangle dt=\chi(\rho)\langle \nabla H_t, \nabla H_t\rangle dt
%	\]
%	and then the proof is complete. 
%	
%\end{proof}
%




 
\bibliographystyle{plain}
\bibliography{zhaoreference.bib}
\end{document}
